\subsection*{19AA. First strict universality-obligation package above \(\mathbf S_4(\mathcal R)\)}
Once the present line has reached the relative universality level \(\mathbf S_4(\mathcal R)\) for a chosen transported tape-based reference core and target class, the remaining question is no longer whether the primitive computational roles are available, nor whether a conservative tape-based reference family can witness relative universality. What remains is the stricter global question whether the transported tape-based route supports the genuinely unbounded extension needed for strict Turing-style completeness. The next conservative step is therefore to isolate that last gap as a separate package instead of mixing it back into the already discharged local hardening ladder.

\begin{definition}[First strict universality-obligation package above \(\mathbf S_4(\mathcal R)\)]
Assume that \(\mathbf S_4(\mathcal R)\) has been reached for an intended target class \(\mathcal K\). Write
\[
\mathbf V=(\mathbf V_\infty,\mathbf V_{\mathrm{ext}},\mathbf V_{\mathrm{dec}})
\]
for the first strict universality-obligation package above \(\mathbf S_4(\mathcal R)\), where:
\begin{itemize}[leftmargin=2em]
  \item \(\mathbf V_\infty\) is the obligation that the chosen transported tape-based reference route admits arbitrarily long admissible extensions rather than only bounded local windows;
  \item \(\mathbf V_{\mathrm{ext}}\) is the obligation that the already established tape-globalization and concatenation discipline extends coherently along such arbitrarily long admissible chains, without hidden global collapse or cutoff;
  \item \(\mathbf V_{\mathrm{dec}}\) is the obligation that the admissible encoding/decoding discipline witnessing \(\mathbf S_4(\mathcal R)\) remains stable under those unbounded extensions and does not depend on a fixed finite window budget.
\end{itemize}
Thus \(\mathbf V\) isolates the first strict machine-theoretic gap between relative universality inside a chosen tape-based reference-core route and full unbounded Turing-style completeness.
\end{definition}

\begin{proposition}[First isolation of the strict universality gap above \(\mathbf S_4(\mathcal R)\)]
Assume that \(\mathbf S_4(\mathcal R)\) has been reached through the conditional reference-core closure surface, the conservative tape-family selection rule, the internal-to-tape transport certificate, and an explicit tape-family relative universality witness. Then the only remaining machine-theoretic step toward the strict level \(\mathbf S_5\) is discharged by the package \(\mathbf V\) above.

Equivalently, once the present line has been upgraded to \(\mathbf S_4(\mathcal R)\), no additional local hardening of control, branch, loop/return, tape support, or reference-core simulation is needed. What remains is exclusively the global unboundedness, extension, and decoding-stability problem recorded by \(\mathbf V\).
\end{proposition}

\begin{corollary}[What the first strict universality-obligation package already buys]
The machine-theoretic hierarchy is now disciplined at both of its remaining upper gaps. The passage from \(\mathbf S_3(\mathcal R)\) to \(\mathbf S_4(\mathcal R)\) is controlled by the conservative tape-family upgrade route, while the passage from \(\mathbf S_4(\mathcal R)\) to \(\mathbf S_5\) is now isolated as the strict universality-obligation package \(\mathbf V\). Later revisions can therefore state very clearly whether they are addressing relative universality inside a chosen transported tape-based reference family or the separate problem of genuine unbounded Turing-style completeness.
\end{corollary}


\subsection*{19AB. Hardening priority for the strict universality obligation above \(\mathbf S_4(\mathcal R)\)}
Once the strict universality gap above \(\mathbf S_4(\mathcal R)\) has been isolated as the package \(\mathbf V=(\mathbf V_\infty,\mathbf V_{\mathrm{ext}},\mathbf V_{\mathrm{dec}})\), the next conservative step is not to attack all three remaining obligations at once. As on the earlier Zuse-guided route, the cleaner strategy is to impose a first hardening order, so that later revisions can reduce the strict universality gap in the same staged way as the primitive computational obligations.

\begin{definition}[First hardening priority order for the strict universality-obligation package above \(\mathbf S_4(\mathcal R)\)]
Assume that the present compressed computational surface has reached \(\mathbf S_4(\mathcal R)\) for an intended target class \(\mathcal K\), and let
\[
\mathbf V=(\mathbf V_\infty,\mathbf V_{\mathrm{ext}},\mathbf V_{\mathrm{dec}})
\]
be the first strict universality-obligation package above \(\mathbf S_4(\mathcal R)\).
Define the first hardening priority order on \(\mathbf V\) by
\[
\mathbf H_2=(\mathbf V_{\mathrm{ext}},\mathbf V_{\mathrm{dec}},\mathbf V_\infty).
\]
Thus the conservative strengthening strategy is:
\begin{itemize}[leftmargin=2em]
  \item first harden the global extension obligation \(\mathbf V_{\mathrm{ext}}\),
  \item then harden decoding stability under extension \(\mathbf V_{\mathrm{dec}}\),
  \item and only at the end attack genuine unboundedness \(\mathbf V_\infty\).
\end{itemize}
\end{definition}

\begin{proposition}[First priority reduction for the strict universality-obligation package above \(\mathbf S_4(\mathcal R)\)]
Assume that \(\mathbf S_4(\mathcal R)\) has been reached and that the hardening order \(\mathbf H_2\) above is adopted. Then later revisions can reduce the strict universality gap monotonically: any successful discharge of the first two obligations \(\mathbf V_{\mathrm{ext}}\) and \(\mathbf V_{\mathrm{dec}}\) shrinks the remaining open strict gap to the single final problem \(\mathbf V_\infty\).

Equivalently, once the extension discipline and its decoding stability have been hardened, the only remaining obstruction to \(\mathbf S_5\) is the genuinely unbounded continuation problem rather than the local or semiglobal coherence of the transported tape-based route.
\end{proposition}

\begin{corollary}[What the first hardening priority order above \(\mathbf S_4(\mathcal R)\) already buys]
The strict universality gap is now not only isolated but also staged. Later revisions can therefore work above \(\mathbf S_4(\mathcal R)\) without mixing extension, decoding, and unboundedness into a single undifferentiated final problem. In the conservative reading, the current next target is no longer \(\mathbf S_5\) as a whole but specifically the first extension-hardening step \(\mathbf V_{\mathrm{ext}}\).
\end{corollary}


\subsection*{19AC. First local hardening of the extension obligation \(\mathbf V_{\mathrm{ext}}\) above \(\mathbf S_4(\mathcal R)\)}
With the strict universality-obligation package now ordered by \(\mathbf H_2\), the next conservative step is to attack the first remaining upper-gap component \(\mathbf V_{\mathrm{ext}}\). The question is no longer whether the present line supports a finite transported tape-based computation relative to a chosen reference core, but whether those admissible windows extend coherently across longer chains without hidden collapse, cutoff, or incompatible seam behavior.

\begin{definition}[First extension-hardening certificate for \(\mathbf V_{\mathrm{ext}}\) above \(\mathbf S_4(\mathcal R)\)]
Assume \(\mathbf S_4(\mathcal R)\) for an intended target class \(\mathcal K\).
Write
\[
\mathfrak E^{\mathrm{glob}}(\mathcal R)=1
\]
if the following hold:
\begin{itemize}[leftmargin=2em]
  \item every admissible finite transported tape-window in the chosen family extends to a longer admissible tape-window by a compatible local continuation step;
  \item the already established concatenation/cut and patch-normal-form discipline remains coherent under those longer continuations;
  \item no hidden global seam defect, cutoff phenomenon, or incompatible local reindexing appears when adjoining admissible windows along the transported tape-line.
\end{itemize}
Thus \(\mathfrak E^{\mathrm{glob}}(\mathcal R)=1\) certifies the first coherent extension-hardening of the tape-based route above \(\mathbf S_4(\mathcal R)\).
\end{definition}

\begin{proposition}[First reduction of the extension obligation \(\mathbf V_{\mathrm{ext}}\) above \(\mathbf S_4(\mathcal R)\)]
Assume \(\mathbf S_4(\mathcal R)\) and \(\mathfrak E^{\mathrm{glob}}(\mathcal R)=1\). Then the first strict universality-obligation package reduces from
\[
\mathbf V=(\mathbf V_\infty,\mathbf V_{\mathrm{ext}},\mathbf V_{\mathrm{dec}})
\]
to
\[
\mathbf V'=(\mathbf V_{\mathrm{dec}},\mathbf V_\infty).
\]
Equivalently, once the transported tape-based route is globally extensible in the admissible-window sense, the strict universality gap above \(\mathbf S_4(\mathcal R)\) no longer contains a separate extension obstruction; what remains are the decoding-stability problem under such extension and the final genuinely unbounded continuation problem.
\end{proposition}

\begin{corollary}[What the first local hardening of \(\mathbf V_{\mathrm{ext}}\) above \(\mathbf S_4(\mathcal R)\) already buys]
The upper strict gap above relative universality is now cleaner than before. If the extension-hardening certificate holds, later revisions do not need to reopen the global seam/coherence problem of transported admissible tape-windows. The next active upper-gap target becomes decoding stability under extension, with genuine unboundedness kept as the final separate endpoint problem.
\end{corollary}


\subsection*{19AD. First local hardening of the decoding-stability obligation \(\mathbf V_{\mathrm{dec}}\) above \(\mathbf S_4(\mathcal R)\)}
Once the extension obligation has been isolated and locally hardened, the next conservative strict-universality target is no longer the seam/coherence of admissible transported tape-windows themselves, but the stability of the admissible encoding/decoding discipline under those longer extensions. The question is now whether the same transported tape-based route that already witnesses relative universality continues to admit a coherent decoding discipline when admissible windows are enlarged, concatenated, and reread along longer chains.

\begin{definition}[First decoding-stability certificate for \(\mathbf V_{\mathrm{dec}}\) above \(\mathbf S_4(\mathcal R)\)]
Assume \(\mathbf S_4(\mathcal R)\) for an intended target class \(\mathcal K\), and assume moreover that the global extension certificate \(\mathfrak E^{\mathrm{glob}}(\mathcal R)=1\) has already been established.
Write
\[
\mathfrak D^{\mathrm{ext}}(\mathcal R)=1
\]
if the following hold:
\begin{itemize}[leftmargin=2em]
  \item the admissible decoding discipline used to witness \(\mathbf S_4(\mathcal R)\) on finite transported tape-windows extends coherently to every admissible longer transported tape-window produced by the extension certificate;
  \item restriction from a longer admissible transported tape-window to any admissible subwindow commutes with the chosen decoding/readout rule;
  \item admissible concatenation of transported tape-windows does not introduce a new decoding ambiguity, hidden branch aliasing, or context-dependent reinterpretation of already decoded local machine states.
\end{itemize}
Thus \(\mathfrak D^{\mathrm{ext}}(\mathcal R)=1\) certifies the first stable decoding discipline under admissible transported tape-extension above \(\mathbf S_4(\mathcal R)\).
\end{definition}

\begin{proposition}[First reduction of the decoding-stability obligation \(\mathbf V_{\mathrm{dec}}\) above \(\mathbf S_4(\mathcal R)\)]
Assume \(\mathbf S_4(\mathcal R)\), \(\mathfrak E^{\mathrm{glob}}(\mathcal R)=1\), and \(\mathfrak D^{\mathrm{ext}}(\mathcal R)=1\). Then the strict universality-obligation package reduces from
\[
\mathbf V=(\mathbf V_\infty,\mathbf V_{\mathrm{ext}},\mathbf V_{\mathrm{dec}})
\]
to
\[
\mathbf V''=(\mathbf V_\infty).
\]
Equivalently, once the transported tape-based route is globally extensible and its admissible decoding discipline remains stable under those extensions, the strict universality gap above \(\mathbf S_4(\mathcal R)\) no longer contains a separate decoding obstruction; what remains is only the final genuinely unbounded continuation problem.
\end{proposition}

\begin{corollary}[What the first local hardening of \(\mathbf V_{\mathrm{dec}}\) above \(\mathbf S_4(\mathcal R)\) already buys]
The upper strict gap has now been reduced as far as it can be reduced without addressing genuine unboundedness itself. If the extension-hardening and decoding-stability certificates both hold, later revisions no longer need to reopen finite-window extension or decoding coherence. The only remaining active upper-gap target is then \(\mathbf V_\infty\), namely true unbounded continuation.
\end{corollary}


\subsection*{19AE. Local hardening of genuine unboundedness \(\mathbf V_\infty\) above \(\mathbf S_4(\mathcal R)\)}
Once global extension and decoding stability under extension have both been isolated and locally hardened, the only remaining strict-universality target is no longer any finite-window transport or rereading problem. What remains is the genuinely unbounded continuation problem itself: whether the transported tape-based route above \(\mathbf S_4(\mathcal R)\) admits arbitrarily extendable admissible continuations without forced collapse, periodic truncation, or hidden resource cutoff.

\begin{definition}[First genuine-unboundedness certificate for \(\mathbf V_\infty\) above \(\mathbf S_4(\mathcal R)\)]
Assume \(\mathbf S_4(\mathcal R)\), \(\mathfrak E^{\mathrm{glob}}(\mathcal R)=1\), and \(\mathfrak D^{\mathrm{ext}}(\mathcal R)=1\).
Write
\[
\mathfrak U^{\infty}(\mathcal R)=1
\]
if the following hold:
\begin{itemize}[leftmargin=2em]
  \item for every admissible transported tape-window in the chosen tape-based family there exists a strictly longer admissible transported tape-window extending it;
  \item admissible extension can be iterated indefinitely without forcing a terminal collapse, bounded cycling artifact, or hidden global cutoff that would prevent arbitrarily long transported computations;
  \item the already established extension and decoding certificates remain valid along every finite stage of such iterated admissible continuation.
\end{itemize}
Thus \(\mathfrak U^{\infty}(\mathcal R)=1\) certifies the first genuinely unbounded admissible continuation regime above \(\mathbf S_4(\mathcal R)\).
\end{definition}

\begin{proposition}[First reduction of the genuine unboundedness obligation \(\mathbf V_\infty\) above \(\mathbf S_4(\mathcal R)\)]
Assume \(\mathbf S_4(\mathcal R)\), \(\mathfrak E^{\mathrm{glob}}(\mathcal R)=1\), \(\mathfrak D^{\mathrm{ext}}(\mathcal R)=1\), and \(\mathfrak U^{\infty}(\mathcal R)=1\). Then the strict universality-obligation package reduces from
\[
\mathbf V=(\mathbf V_\infty,\mathbf V_{\mathrm{ext}},\mathbf V_{\mathrm{dec}})
\]
to the empty list.
Equivalently, once the transported tape-based route is globally extensible, decoding-stable under extension, and genuinely unbounded in the admissible continuation sense, the strict universality gap above \(\mathbf S_4(\mathcal R)\) contains no further separate obstruction. In that case the compressed computational surface reaches the top level \(\mathbf S_5\), namely strict unbounded Turing-style completeness relative to the certified route.
\end{proposition}

\begin{corollary}[What the first local hardening of \(\mathbf V_\infty\) above \(\mathbf S_4(\mathcal R)\) already buys]
The upper gap is now completely structured. If the global extension, decoding-stability, and genuine-unboundedness certificates all hold, later revisions do not need to reopen any remaining strict-universality obstruction above \(\mathbf S_4(\mathcal R)\). The only remaining caution is interpretive: one must still keep explicit which route, tape-family, and reference-core class are being certified when speaking of strict completeness.
\end{corollary}


\subsection*{19AF. Full machine-theoretic closure surface on the compressed computational line}
After the conditional reference-core closure surface, the relative universality-upgrade witness, and the strict upper certificate package have all been isolated, the machine-theoretic reading should no longer appear as a long ladder of locally discharged obligations. It can now be bundled into a single closure surface that records exactly which level of the hierarchy has been reached under which certificate discipline.

\begin{definition}[Full machine-theoretic closure surface on the compressed computational line]
Fix a chosen admissible sequential reference core \(\mathcal R\), a chosen admissible tape-based reference family, and the associated conservative transport discipline. Write
\[
\mathfrak M^{\mathrm{full}}(\mathcal R)=1
\]
if the following package holds simultaneously:
\begin{itemize}[leftmargin=2em]
  \item the compressed computational interpretation surface reaches the conditional reference-core closure level \(\mathbf S_3(\mathcal R)\), i.e. \(\mathfrak Z^{\mathrm{ref}}(\mathcal R)=1\);
  \item the tape-family upgrade route is certified, i.e. the conservative tape-family selection rule, the internal-to-tape transport certificate, and the tape-family relative universality witness all hold, so that the line reaches \(\mathbf S_4(\mathcal R)\);
  \item the strict upper certificate package is certified, i.e.
  \[
  \mathfrak E^{\mathrm{glob}}(\mathcal R)=1,\qquad
  \mathfrak D^{\mathrm{ext}}(\mathcal R)=1,\qquad
  \mathfrak U^{\infty}(\mathcal R)=1,
  \]
  so that the remaining strict gap above \(\mathbf S_4(\mathcal R)\) is empty.
\end{itemize}
Thus \(\mathfrak M^{\mathrm{full}}(\mathcal R)=1\) denotes the first fully bundled machine-theoretic closure surface attached to the present compressed computational line relative to the chosen certified route.
\end{definition}

\begin{theorem}[First bundled machine-theoretic closure theorem on the compressed computational line]
Assume \(\mathfrak M^{\mathrm{full}}(\mathcal R)=1\). Then the compressed computational line reaches the top machine-theoretic level \(\mathbf S_5\) along the chosen certified route. Equivalently, all intermediate gaps
\[
\mathbf S_1\longrightarrow \mathbf S_2\longrightarrow \mathbf S_3(\mathcal R)\longrightarrow \mathbf S_4(\mathcal R)\longrightarrow \mathbf S_5
\]
are discharged under an explicitly named certificate package, so that no separate machine-theoretic obstruction remains inside the adopted route.
\end{theorem}

\begin{corollary}[What the full machine-theoretic closure surface already buys]
Later revisions no longer need to cite the whole local hardening ladder when the only point is that the present computational reading has been conditionally closed all the way up the hierarchy along a chosen certified route. They may instead cite \(\mathfrak M^{\mathrm{full}}(\mathcal R)=1\), while still keeping explicit which reference core, which tape-family route, and which upper certificates are meant. This compresses the machine-theoretic reading to a single surface without pretending that the route has become unconditional or model-independent.
\end{corollary}




\subsection*{19AG. First machine-theoretic status decoder below and above the full closure surface}
After the full machine-theoretic closure surface has been bundled, the text should not force later revisions to reopen the entire hierarchy merely to say which machine-theoretic level is currently justified under a given certificate package. A small status decoder is therefore useful: it reads off the strongest currently justified strength level from the available closure surfaces and upgrade certificates without changing the underlying mathematics.

\begin{definition}[First machine-theoretic status decoder on the compressed computational line]
Fix a chosen admissible sequential reference core \(\mathcal R\). Define the machine-theoretic status decoder
\[
\mathfrak S^{\mathrm{mach}}(\mathcal R)\in\{\mathbf S_1,\mathbf S_2,\mathbf S_3(\mathcal R),\mathbf S_4(\mathcal R),\mathbf S_5\}
\]
as the strongest level among the hierarchy
\[
\mathbf S_1\longrightarrow \mathbf S_2\longrightarrow \mathbf S_3(\mathcal R)\longrightarrow \mathbf S_4(\mathcal R)\longrightarrow \mathbf S_5
\]
whose required certificate package has been verified. Concretely:
\begin{itemize}[leftmargin=2em]
  \item \(\mathfrak S^{\mathrm{mach}}(\mathcal R)=\mathbf S_5\) if \(\mathfrak M^{\mathrm{full}}(\mathcal R)=1\);
  \item \(\mathfrak S^{\mathrm{mach}}(\mathcal R)=\mathbf S_4(\mathcal R)\) if the tape-family upgrade route is certified but the strict upper certificate package is not fully certified;
  \item \(\mathfrak S^{\mathrm{mach}}(\mathcal R)=\mathbf S_3(\mathcal R)\) if the conditional reference-core closure surface is certified but the tape-family upgrade route is not fully certified;
  \item \(\mathfrak S^{\mathrm{mach}}(\mathcal R)=\mathbf S_2\) if the conditional primitive Zuse-style core is certified but reference-core closure is not;
  \item \(\mathfrak S^{\mathrm{mach}}(\mathcal R)=\mathbf S_1\) otherwise.
\end{itemize}
\end{definition}

\begin{proposition}[Monotone exactness of the first machine-theoretic status decoder]
For every chosen admissible sequential reference core \(\mathcal R\), the decoder \(\mathfrak S^{\mathrm{mach}}(\mathcal R)\) is monotone with respect to certificate accumulation and exact with respect to the present hierarchy: it never outputs a higher level than has actually been certified, and whenever a full certificate package for some level in the hierarchy has been established, the decoder outputs at least that level. In particular,
\[
\mathfrak M^{\mathrm{full}}(\mathcal R)=1
\quad\Longrightarrow\quad
\mathfrak S^{\mathrm{mach}}(\mathcal R)=\mathbf S_5.
\]
Likewise, if only the reference-core closure surface is available, the decoder halts at \(\mathbf S_3(\mathcal R)\), and if the tape-family upgrade route has been certified but the strict upper package has not, it halts at \(\mathbf S_4(\mathcal R)\).
\end{proposition}

\begin{corollary}[What the first machine-theoretic status decoder already buys]
Later revisions can now state the strongest presently justified machine-theoretic claim without reopening the whole ladder of lower certificates. They may cite the single decoded status \(\mathfrak S^{\mathrm{mach}}(\mathcal R)\) and then only unfold the route-specific certificate package when a stronger or more unconditional claim is intended. This keeps the computational reading disciplined: the text can speak precisely about the current level already reached, while preventing accidental overstatement about higher levels that are not yet certified.

\end{corollary}


\subsection*{19AH. Conditional/unconditional status surface on the compressed computational line}
After the machine-theoretic status decoder has been introduced, the current release line should not remain described only in abstract terms. A small present-line status surface is therefore useful: it records, in one place, the strongest route-dependent status already certified on the present v\docversion\ release line and separates it from the weaker certificate-free status that can be asserted without invoking any chosen reference-core route.

\begin{definition}[First present-line conditional/unconditional status surface on the compressed computational line]
Fix a chosen admissible sequential reference core \(\mathcal R\). Define the present-line machine-theoretic status pair
\[
\mathfrak P^{\mathrm{mach}}(\mathcal R)=\bigl(\mathfrak S^{\mathrm{abs}},\mathfrak S^{\mathrm{cond}}(\mathcal R)\bigr)
\]
by the following rule:
\begin{itemize}[leftmargin=2em]
  \item \(\mathfrak S^{\mathrm{abs}}\) is the strongest machine-theoretic level justified without invoking any chosen reference core, tape-family route, or upper certificate package;
  \item \(\mathfrak S^{\mathrm{cond}}(\mathcal R)\) is the strongest machine-theoretic level justified under the explicit chosen route \(\mathcal R\) together with the corresponding lower and upper certificate package.
\end{itemize}
The pair \(\mathfrak P^{\mathrm{mach}}(\mathcal R)\) will be called the present-line conditional/unconditional status surface.
\end{definition}

\begin{proposition}[First explicit present-line status readout on the compressed computational line]
On the present v\docversion\ release line, and for every chosen admissible sequential reference core \(\mathcal R\) along which the bundled closure surface \(\mathfrak M^{\mathrm{full}}(\mathcal R)=1\) has been certified, the present-line status surface reads
\[
\mathfrak P^{\mathrm{mach}}(\mathcal R)=\bigl(\mathbf S_1,\mathbf S_5\bigr).
\]
Equivalently: unconditionally, the present line should still only be advertised at the baseline disciplined computational-reading level \(\mathbf S_1\); conditionally, along the explicitly certified route \(\mathcal R\), the present line reaches the top certified level \(\mathbf S_5\). In particular, the present text should avoid collapsing these two readings into a single unconditional global claim.
\end{proposition}

\begin{corollary}[What the first present-line conditional/unconditional status surface already buys]
Later revisions can now cite the present machine-theoretic state of the Companion in one disciplined sentence: the line is unconditionally at \(\mathbf S_1\), while conditionally along a chosen certified route it reaches \(\mathbf S_5\). This gives a compact status formula for future hardening passes while preventing accidental overstatement about certificate-dependent closures as if they were already unconditional.
\end{corollary}

\begin{remark}[One-paragraph synopsis of the machine-theoretic line]
The computational line of the present framework can be read in three nested stages. First, the local operational layer yields a conditional primitive machine core, in which tape-like support, control, branching, and loop/return behavior are structurally separated and locally executable. Second, this primitive core can be hardened into a reference-core closure relative to a chosen sequential reference kernel \(\mathcal R\), so that the line is no longer only machine-like in analogy, but carries a disciplined reference computation surface. Third, along a certified tape-based route above this reference-core closure, the framework admits a conditional upgrade toward stronger machine-theoretic levels, up to a Turing-style completeness claim only under explicit higher certificates; accordingly, the present line supports an unconditional computational base reading, while stronger universality claims remain route-dependent and conditional.
\end{remark}




\subsection*{19AH$'$. First bidirectional observer-side machine-shadow split below the status surface}
After the present-line conditional/unconditional status surface has been fixed, one more local question becomes unavoidable on the working continuation line: how should the currently retained public machine readout itself be decomposed at observer side? The next useful step is not yet a full new universality statement, but a disciplined split of the retained machine shadow into one symmetric total channel and one antisymmetric imbalance channel. This keeps the older no-second-carrier discipline visible while making the current machine readout auditable in a form that can be compared directly with later transport, spectral, or shell-level hardening.

\begin{definition}[Bidirectional observer-side machine-shadow coordinates]
Fix the retained overlap window \(W_{\mathrm{ret}}\) and a tested lag \(\ell\). Let \(\eta_{\mathrm{drive}}^{(\ell)}\) denote the fitted observer-side drive contribution and \(\eta_{\mathrm{carrier}}^{(\ell)}\) the fitted carrier contribution on \(W_{\mathrm{ret}}\). Define the bidirectional observer-side coordinates
\[
\eta_+^{(\ell)}=\eta_{\mathrm{drive}}^{(\ell)}+\eta_{\mathrm{carrier}}^{(\ell)},
\qquad
\eta_-^{(\ell)}=\eta_{\mathrm{drive}}^{(\ell)}-\eta_{\mathrm{carrier}}^{(\ell)}.
\]
We call \(\eta_+^{(\ell)}\) the \emph{total machine-shadow channel} and \(\eta_-^{(\ell)}\) the \emph{observer-side imbalance channel}.
\end{definition}

\begin{proposition}[Current retained-overlap checkpoint at the public optimum]
On the retained overlap window currently kept on the working line, the bidirectional observer-side coordinates reproduce the present public optimum at lag
\[
\ell_*=-14
\]
with weighted defect approximately \(0.6746\) over \(42\) overlap points.
\end{proposition}

\begin{remark}[Interpretation of the checkpoint]
The point of this checkpoint is not yet a final dynamical theorem about the origin of the machine readout. It is the weaker but already useful fact that the presently best public lag can be read consistently through one bidirectional observer-side split rather than through two unrelated fitted channels. In that sense, the present step converts one numerical optimum into a structurally readable local machine-shadow record.
\end{remark}

\begin{proposition}[The imbalance channel is substantial but not primitive]
On the same retained overlap window at the public optimum \(\ell_*=-14\), the observer-side imbalance channel satisfies
\[
\frac{\lVert \eta_-^{(\ell_*)}\rVert_1}{\lVert \eta_+^{(\ell_*)}\rVert_1}\approx 0.6827,
\qquad
\operatorname{corr}\!\bigl(\eta_+^{(\ell_*)},\eta_-^{(\ell_*)}\bigr)\approx 0.9532.
\]
Hence the imbalance channel is not negligible, but it remains tightly slaved to the same total machine-shadow profile.
\end{proposition}

\begin{corollary}[First no-second-primitive-carrier reading from the machine split]
The current bidirectional observer-side readout does not force the introduction of a second primitive carrier. What it supports is one dominant machine shadow together with a substantial but highly correlated imbalance readout. In particular, the current continuation line remains compatible with the older no-second-carrier discipline: the extra visible degree of freedom appears at observer side as a locked asymmetry channel, not as evidence for a new independent carrier species.
\end{corollary}


\begin{remark}[Why this matters for the next hardening step]
The split creates a useful next leverage point. Future hardening no longer needs to ask only whether a retained carrier exists; it can ask whether a proposed theorem, transport step, or spectral proxy controls the symmetric total channel \(\eta_+\), the locked imbalance channel \(\eta_-\), or both at once. That distinction should make it easier to separate genuine carrier obligations from observer-shell correction obligations in the next working passes.
\end{remark}

\subsection*{19AH$''$. First channel-control dictionary induced by the machine-shadow split}
Once the observer-side machine readout has been split into \(\eta_+\) and \(\eta_-\), the next useful step is a control dictionary telling later arguments what kind of burden they actually carry. This is still weaker than a new dynamical theorem. Its role is bookkeeping with mathematical consequences: it distinguishes propositions that genuinely control the retained machine shadow from propositions that merely reorganize the observer-side asymmetry readout around that shadow.

\begin{definition}[\(\eta_+\)-control, \(\eta_-\)-control, and joint channel control]
Fix the retained overlap window \(W_{\mathrm{ret}}\) at the current public optimum \(\ell_*=-14\). We say that a theorem, estimate, transport rule, or spectral proxy on the present working line has
\begin{enumerate}[leftmargin=2em]
  \item \emph{total-channel control} if it bounds, reconstructs, or rigidifies \(\eta_+^{(\ell_*)}\) up to the standing admissible equivalences without requiring an independent primitive contribution beyond the present drive/carrier split;
  \item \emph{imbalance-channel control} if, once the total channel is fixed in that admissible sense, it bounds, reconstructs, or rigidifies \(\eta_-^{(\ell_*)}\) as an observer-side asymmetry readout around the same total machine shadow;
  \item \emph{joint channel control} if it controls both channels at once and in addition preserves or explains a structural lock between them, such as the present norm ratio or the present high correlation.
\end{enumerate}
\end{definition}

\begin{proposition}[First machine-side reading discipline induced by the dictionary]
On the current working line, arguments aimed at the retained carrier or at the existence of one common machine shadow should be read as \(\eta_+\)-control problems, whereas observer-shell, asymmetry, phase-correction, or frame-flicker refinements should be read as \(\eta_-\)-control problems unless they also alter the total channel itself. Mixed claims that try to upgrade both at once should be treated as joint-channel problems and therefore carry the heaviest burden of proof.
\end{proposition}

\begin{corollary}[First reformulation of the no-second-carrier discipline at machine level]
A future refinement is compatible with the present no-second-carrier discipline whenever its extra visible freedom can be organized as \(\eta_-\)-control over a fixed or jointly controlled \(\eta_+\)-channel, rather than by introducing a second primitive contribution already at the total-channel level. In particular, the present checkpoint pushes the burden of any would-be second-carrier claim upward: such a claim would now have to break the current reading at the level of \(\eta_+\), not merely by producing additional observer-side imbalance structure.
\end{corollary}

\begin{remark}[Why the dictionary is a real compression tool]
The dictionary turns one vague question --- ``does this next theorem improve the machine line?'' --- into a sharper trichotomy. A proposed step may improve the total machine shadow, the locked imbalance readout, or the coupling data between them. This should make later proof compression cleaner because many older optical and shell-side corrections can now be classified immediately as \(\eta_-\)-side obligations, while the rarer genuinely carrier-bearing steps can be isolated as \(\eta_+\)-side obligations.
\end{remark}

\subsection*{19AH$'''$. First backward channel reading of the older optical and shell-side baselines}
The dictionary becomes substantially more useful once it reorganizes not only future arguments but also the already stabilized comparison/shell archive. The present working line therefore records a first backward reading of the older optical-comparison, shortlist, bombe/sieve, shell/menu/window, and flicker-style material through the split into \(\eta_+\) and \(\eta_-\).

\begin{proposition}[First backward channel reading of the retained public baselines]
In this inherited public-release baseline block, the older public baselines should be read according to the following burden split.
\begin{enumerate}[leftmargin=2em]
  \item The retained-carrier comparison layer, public optimum selection, no-second-carrier discipline, and every comparison/back-reading/bombe/sieve step insofar as it isolates one common machine-bearing shadow are primarily \(\eta_+\)-control statements.
  \item The shell/menu/window refinements, phase-correction re-readings, and frame-flicker-style observer corrections are primarily \(\eta_-\)-control statements whenever they refine how that common shadow appears on the observer side without changing the retained optimum itself.
  \item Any statement claiming at once that a shell-side correction preserves the retained optimum, explains the asymmetry pattern, and leaves the common machine shadow identifiable should be treated as a joint-channel statement.
\end{enumerate}
\end{proposition}

\begin{corollary}[First re-reading of the optical approximation baseline]
The earlier first retained-carrier optical approximation should presently be read as one \(\eta_+\)-bearing baseline together with a subordinate \(\eta_-\)-side correction envelope. Accordingly, visible mismatch inside that optical approximation is not yet evidence against the retained-carrier discipline unless it forces an additional primitive contribution already at the total-channel level.
\end{corollary}

\begin{remark}[Compression payoff of the backward reading]
The backward reading turns a long mixed archive into a more ordered burden ledger. Carrier isolation, common-window retention, and no-second-carrier arguments can now be cited as total-channel obligations, while shell refinements, asymmetry diagnostics, and flicker-style observer corrections can be cited as imbalance-side obligations. Only the rarer steps that claim both retention and corrected observer-side readout need to reopen the full joint burden.
\end{remark}


\subsection*{19AH$^{(4)}$. First theorematic burden split for the remaining continuation knots}
Once the older baselines have been re-read through \(\eta_+\) and \(\eta_-\), the next practical question is which unresolved continuation knots still carry genuine carrier burden and which ones are only observer-side refinement burdens. The present working line therefore records one first theorematic burden split: unresolved shell-, phase-, or flicker-side refinement is not automatically a retained-carrier obstruction.

\begin{proposition}[First theorematic burden split for the remaining continuation knots]
On the present continuation line, the remaining theorematic knots should be read according to the following burden discipline.
\begin{enumerate}[leftmargin=2em]
  \item Any future statement whose conclusion asserts preservation of the retained public optimum, one-common-machine-shadow identification, or the no-second-primitive-carrier discipline must control the total channel \(\eta_+\). Purely imbalance-side refinements on \(\eta_-\) can neither prove nor refute such a carrier-level statement by themselves.
  \item The frame-flicker ansatz, shell/menu/window refinements, phase-correction re-readings, and analogous observer-side asymmetry refinements are to be treated as \(\eta_-\)-side knots as long as they only change how the retained machine shadow is seen and do not force a new total-channel profile.
  \item Any theorem claiming simultaneously that an observer-side correction explains the asymmetry pattern, improves the retained fit, and preserves one common carrier must be treated as a joint-channel knot: it has to control both \(\eta_+\) and \(\eta_-\).
  \item Consequently, a visible residual defect that can be confined to \(\eta_-\) is not yet a retained-carrier failure of the present line. A genuine carrier-level obstruction appears only if the retained optimum, the one-shadow reading, or the no-second-carrier discipline is already forced to fail on the total-channel side.
\end{enumerate}
\end{proposition}

\begin{corollary}[Default downgrade rule for unresolved observer-side mismatch]
On the present continuation line, an unresolved optical, shell-side, or flicker-style mismatch should by default be downgraded to an \(\eta_-\)-side obligation unless it demonstrably changes the retained optimum, destroys the one-shadow reading at the total-channel level, or forces an additional primitive contribution already in \(\eta_+\). This keeps the current no-second-carrier discipline from being overburdened by observer-side correction tasks that have not yet crossed into carrier space.
\end{corollary}

\begin{remark}[Why this sharpened burden split matters]
The practical gain is a much cleaner proof queue. Future hardening can now ask first whether an open point is really about the common retained shadow or only about its observer-side distortion. If it is only the latter, then it belongs to the asymmetry shell and should not delay carrier-level statements. Only the comparatively rare steps that touch both channels at once need the full coupled burden again.
\end{remark}

\subsection*{19AH$^{(5)}$. First formal demotion criterion from joint-channel burden to pure imbalance-side burden}
The burden split of 19AH$^{(4)}$ becomes operational only once one can say when a candidate continuation step no longer deserves to be treated as a genuinely joint-channel knot. The next local tool is therefore a formal demotion criterion: a pending observer-side refinement may be demoted from joint-channel burden to a pure \(\eta_-\)-side burden as soon as its unresolved remainder provably leaves the retained total profile untouched.

\begin{definition}[Observer-side demotion certificate]
Fix a pending continuation statement \(\mathcal C\) whose unresolved part currently concerns shell, phase, flicker, menu, or analogous observer-side refinement. We say that \(\mathcal C\) admits an \emph{observer-side demotion certificate} if all three of the following conditions hold on the retained admissible window under discussion.
\begin{enumerate}[leftmargin=2em]
  \item \textbf{Total-profile invariance.} The retained public optimum, the one-machine-shadow identification, and the no-second-primitive-carrier reading remain unchanged on the \(\eta_+\)-side under the unresolved remainder of \(\mathcal C\).
  \item \textbf{No forced new primitive contribution.} The unresolved remainder of \(\mathcal C\) does not require an additional primitive carrier already at the total-channel level; equivalently, the outstanding defect can still be represented without altering the retained \(\eta_+\)-carrier architecture.
  \item \textbf{Imbalance localization.} The visible mismatch produced by the unresolved remainder of \(\mathcal C\) can be localized to the observer-side imbalance channel \(\eta_-\), meaning that its effect is read as a change of asymmetry, shell, or observer-side correction profile rather than as a change of the retained common total shadow.
\end{enumerate}
\end{definition}

\begin{proposition}[First formal demotion criterion]
If a pending continuation statement \(\mathcal C\) admits an observer-side demotion certificate, then \(\mathcal C\) may be demoted from a provisional joint-channel knot to a pure \(\eta_-\)-side obligation on the present continuation line. In particular, the unresolved remainder of \(\mathcal C\) may no longer be cited as an open carrier-level obstruction against the retained public optimum, the one-shadow reading, or the no-second-primitive-carrier discipline.
\end{proposition}

\begin{proof}
The three certificate clauses remove exactly the reasons why \(\mathcal C\) would still need joint status. Clause~(1) says that the retained total profile already survives unchanged at the \(\eta_+\)-level. Clause~(2) rules out the only remaining carrier-level escape route, namely that the unresolved part secretly forces an additional primitive contribution already in the total channel. Clause~(3) places the visible remainder entirely on the observer-side imbalance ledger. Once these three facts are available simultaneously, the outstanding burden of \(\mathcal C\) no longer concerns whether the present line carries one retained common machine shadow, but only how the observer-side asymmetry shell should be refined. The knot is therefore genuinely \(\eta_-\)-side and may be demoted accordingly.
\end{proof}

\begin{corollary}[Safe demotion rule for future continuation work]
Future continuation steps should only be kept on the joint-channel ledger if at least one of the three certificate clauses fails or is not yet known. Conversely, once the retained \(\eta_+\)-profile is stable, no extra primitive carrier is forced there, and the remaining defect is localized to \(\eta_-\), the step should be queued as observer-side hardening rather than as unfinished carrier proof.
\end{corollary}

\begin{remark}[Why this tool is useful]
This criterion turns the burden split into a reusable proof tool. It prevents later drafts from repeatedly reopening the full carrier question whenever a shell-side mismatch or flicker-style correction is still under discussion. The proof queue can now be triaged by certificate: unresolved points without \(\eta_+\)-damage stay on the asymmetry ledger, while only genuine total-channel risk keeps a knot on the joint ledger.
\end{remark}


\subsection*{19AH$^{(6)}$. First provisional demotion ledger for inherited observer-side continuation knots}
\noindent\textbf{Reader cue.} Read the next two subsections as a queue-classification pair. Their purpose is to sort inherited shell/phase/flicker-style leftovers first onto the observer-side ledger by default and then to state the exact evidence that would force those same knots back into joint or total-channel burden.

The formal demotion criterion of 19AH$^{(5)}$ becomes most useful once it is applied back to the inherited queue. The point is not to declare old open points ``solved'' by fiat, but to state explicitly which of them should from now on be carried on the observer-side ledger unless and until they actually damage the retained total profile. The next local step is therefore a provisional demotion ledger for the main inherited observer-side knot families.

\begin{definition}[Provisional demotion ledger for inherited observer-side knots]
On the present continuation line, the following inherited knot families are placed on the \emph{provisional observer-side demotion ledger}, subject to re-promotion only if later work shows a genuine failure of at least one certificate clause from 19AH$^{(5)}$:
\begin{enumerate}[leftmargin=2em]
  \item shell/menu/window-refinement mismatches that alter the observer-side shell geometry without changing the retained total optimum profile,
  \item phase-correction and frame-flicker residuals that modify how the retained profile is read from the observer side while leaving the common total shadow intact,
  \item optical-approximation residuals whose visible discrepancy can still be localized to the asymmetry envelope rather than to the retained common carrier, and
  \item analogous observer-side refinements inherited from the older archive whenever they do not force a second primitive carrier and do not move the retained public optimum on the total-channel side.
\end{enumerate}
\end{definition}

\begin{proposition}[First provisional demotion ledger theorem]
The inherited knot families listed above belong, by default, on the observer-side demotion ledger rather than on the joint-channel obstruction ledger. More precisely, on the present compressed computational line they should be read as provisional \(\eta_-\)-side obligations unless later analysis shows that one of them changes the retained \(\eta_+\)-profile, breaks the one-shadow reading at total-channel level, or forces an additional primitive carrier already in \(\eta_+\).
\end{proposition}

\begin{proof}
Each listed family is inherited precisely from the shell/phase/flicker/optical side of the archive, hence from the side that 19AH$'''$ and 19AH$^{(4)}$ already identified as the natural observer-side burden domain. The formal criterion of 19AH$^{(5)}$ then supplies the missing operational rule: as long as the retained total profile remains fixed, no new primitive carrier is forced in the total channel, and the visible mismatch stays localizable to the imbalance channel, the corresponding inherited knot is not entitled to remain on the joint ledger. The listed families therefore enter the queue under default demotion. Re-promotion is reserved only for the exceptional case that later work actually detects \(\eta_+\)-damage or a genuine carrier-level escape route.
\end{proof}

\begin{corollary}[How to read the old queue from now on]
From this point onward, older shell/menu/window, phase-correction, flicker-style, and observer-side optical residuals should no longer be cited as open carrier objections by default. They remain real work items, but they are to be read first as asymmetry-hardening tasks. Only after a concrete failure of the demotion certificate has been exhibited should one move such a knot back to the coupled or carrier-level queue.
\end{corollary}

\begin{remark}[What ``demotion'' means here in plain language]
This is only a bookkeeping and proof-burden move. ``Demotion'' does \emph{not} mean that the point is unimportant, false, or removed. It means: this point still matters, but at the current stage it counts as a problem of observer-side correction rather than as evidence that the common retained carrier picture has broken. In ordinary language, we stop treating every leftover wobble in the readout shell as if it were a collapse of the whole machine-bearing line.
\end{remark}

\subsection*{19AH$^{(7)}$. First formal re-promotion criterion from observer-side demotion back to joint or total burden}
Once a provisional demotion ledger exists, the opposite discipline becomes equally important. Demotion is safe only if later work can also say when a demoted knot must be upgraded again. The next local tool is therefore a formal re-promotion criterion: if a supposedly observer-side remainder turns out to damage the retained total machine-shadow profile or to force a new primitive carrier already at total-channel level, then the knot must leave the pure \(\eta_-\)-ledger and return to a higher burden class.

\begin{definition}[Observer-side re-promotion trigger]
Let a continuation knot family \(K\) currently lie on the provisional observer-side demotion ledger. We say that \(K\) carries an \emph{observer-side re-promotion trigger} if later analysis exhibits at least one of the following three failures:
\begin{enumerate}[leftmargin=2em]
  \item the retained total machine-shadow profile \(\eta_+\) is no longer fixed up to the standing admissible equivalences;
  \item the one-shadow reading at total-channel level breaks, in the sense that the visible fit can no longer be organized without introducing an additional primitive carrier already in \(\eta_+\);
  \item the observer-side remainder can no longer be localized purely to \(\eta_-\), because its correction necessarily changes the total-channel optimum, the total-channel rigidity data, or the total-channel carrier count.
\end{enumerate}
\end{definition}

\begin{proposition}[First re-promotion criterion]
Let \(K\) be a knot family on the provisional observer-side demotion ledger. If \(K\) satisfies an observer-side re-promotion trigger, then \(K\) is no longer admissible as a pure \(\eta_-\)-side obligation. More precisely:
\begin{enumerate}[leftmargin=2em]
  \item if the trigger only shows that \(K\) cannot be separated from the total-channel data but does not yet force a new primitive carrier, then \(K\) must be re-promoted from pure \(\eta_-\)-burden to \emph{joint-channel burden};
  \item if the trigger forces a change in the retained total machine-shadow profile or compels an additional primitive carrier already in \(\eta_+\), then \(K\) must be re-promoted all the way to \emph{total-channel burden}.
\end{enumerate}
\end{proposition}

\begin{proof}
This is the exact converse of the demotion logic from 19AH$^{(5)}$ and the ledger rule from 19AH$^{(6)}$. A knot may remain on the pure observer-side ledger only while the retained total profile stays fixed, no second primitive carrier is forced at total-channel level, and the visible remainder remains localizable to the imbalance channel. The three trigger clauses record the failure of precisely those permissions. Once a clause fails, pure \(\eta_-\)-treatment is no longer justified. If the failure merely entangles the remainder with the total channel, the knot becomes joint. If the failure genuinely alters or multiplies the total-channel carrier structure, the knot becomes a total-channel problem.
\end{proof}

\begin{corollary}[Plain-language queue rule]
A leftover shell, phase, flicker, or optical residual may stay in the ``observer-side corrections'' queue only while it leaves the main common machine shadow intact. The moment it starts moving that main shadow, or starts behaving as if a second primitive carrier were needed in the total profile, it must be upgraded again. In ordinary language: demotion is reversible, but only by evidence of real damage to the main channel.
\end{corollary}

\begin{remark}[Why the demotion/re-promotion pair is useful]
Together, 19AH$^{(5)}$--19AH$^{(8)}$ now give a two-way burden discipline. Demotion prevents every observer-side wobble from being misread as a carrier collapse; re-promotion prevents the opposite mistake of hiding a true carrier-level defect inside the observer-side queue. This makes the inherited continuation landscape auditable in both directions.
\end{remark}

\subsection*{19AH$^{(8)}$. First channel traffic-light ledger induced by the demotion/re-promotion discipline}
\noindent\textbf{Reader cue.} Read the next ledger as a dashboard. Its purpose is to summarize the already stated demotion and re-promotion rules in green/yellow/red form, so that the inherited queue can be scanned quickly before later live burdens are isolated.

The next local consolidation step is to turn the two-way burden discipline into a directly readable queue picture. Instead of carrying every inherited continuation knot under one undifferentiated status word, one may sort the inherited families by a traffic-light rule. Green means that the family presently remains admissible as a pure observer-side burden; yellow means that the family is already entangled enough with the total channel that it must be treated jointly; red is reserved for the genuinely exceptional case that the family damages the retained total profile itself or forces a second primitive carrier already in \(\eta_+\).

\begin{definition}[Channel traffic-light ledger]
A \emph{channel traffic-light ledger} for inherited continuation knots is an assignment of each inherited knot family \(K\) to exactly one of the following three burden classes:
\begin{enumerate}[leftmargin=2em]
  \item \textbf{green / observer-side only}: \(K\) satisfies the demotion certificate from 19AH$^{(5)}$ and no re-promotion trigger from 19AH$^{(7)}$ is presently visible;
  \item \textbf{yellow / joint-channel}: \(K\) no longer qualifies for pure observer-side treatment, but the available evidence still stops short of forcing a genuine total-channel carrier change;
  \item \textbf{red / total-channel}: \(K\) damages the retained total profile \(\eta_+\) itself or forces an additional primitive carrier already at total-channel level.
\end{enumerate}
\end{definition}

\begin{proposition}[First provisional traffic-light placement for the inherited queue]
On the present compressed computational line, the inherited queue may provisionally be sorted as follows.
\begin{enumerate}[leftmargin=2em]
  \item \textbf{Green:} shell/menu/window refinements, phase-correction refinements, frame-flicker compensation refinements, and observer-side optical residual refinements remain on the green ledger, because they presently leave the retained total machine-shadow profile fixed, force no second primitive carrier in \(\eta_+\), and stay localizable to imbalance-side repair.
  \item \textbf{Yellow:} any inherited knot whose proposed correction already shifts the retained optimum location, modifies the rigidity data needed to keep the one-shadow total reading stable, or otherwise couples the observer-side mismatch back into the total-channel fit must enter the yellow ledger and be treated as a joint-channel obligation.
  \item \textbf{Red:} no presently inherited knot family is forced onto the red ledger on the current line; red is reserved only for the case that later work exhibits an actual breakdown of the retained total profile or a genuine second primitive carrier already in \(\eta_+\).
\end{enumerate}
\end{proposition}

\begin{proof}
The green placement is exactly the provisional demotion ledger from 19AH$^{(6)}$, now read through the demotion certificate of 19AH$^{(5)}$. The yellow placement is the first nontrivial output of the re-promotion rule from 19AH$^{(7)}$: once a knot cannot be repaired purely on the imbalance side, yet still does not force a total-channel carrier split, it must be carried jointly. The red placement is currently empty because the present line has not yet produced evidence that the retained total machine shadow itself fails or that a second primitive carrier is already required in the total channel. Hence the inherited queue is presently green-dominant, with yellow reserved for future entanglement cases and red reserved for genuine carrier-level failure.
\end{proof}

\begin{corollary}[Plain-language traffic-light reading]
At the moment, most inherited leftovers are green: they are still real work, but only on the observer side. Yellow means: this leftover is beginning to touch the main common shadow and therefore needs joint treatment. Red would mean: the main common shadow itself is no longer enough. On the current line, we are not at red.
\end{corollary}

\begin{remark}[Why this corrects possible drift]
The point of the traffic-light ledger is precisely to stop the continuation block from drifting into an ever longer meta-language. Instead of adding more abstract burden labels, the present line now has a single concrete reading device: green for pure \(\eta_-\)-repair, yellow for genuine coupling, red for actual \(\eta_+\)-damage. This keeps the new diagnostic layer operational and tied to concrete queue management.
\end{remark}

\subsection*{19AH$^{(9)}$. First synchronized OP ledger after the channel-discipline consolidation}
\noindent\textbf{Reader cue.} Read the next synchronized OP ledger as a theorem-summary device. Its purpose is to collapse historical OP labels into \emph{closed}, \emph{absorbed}, and \emph{residual-core} status before the roadmap is read, not to reopen already discharged fronts.

Once the inherited continuation queue has been sorted by total/imbalance burden and by the traffic-light rule, the older named OP fronts should no longer be carried unchanged merely for historical convenience. The next consolidation step is therefore to synchronize the inherited OP labels against the presently certified structure itself. The point is not to rename difficulty away, but to state explicitly which earlier OP fronts are already structurally discharged, which no longer survive as independent fronts, and which remain as the genuine residual core for the next version.

\begin{definition}[Synchronized OP ledger]
A \emph{synchronized OP ledger} on the present compressed computational line assigns every inherited named OP front \(P\) to exactly one of the following three statuses:
\begin{enumerate}[leftmargin=2em]
  \item \textbf{closed}: the relevant structural burden of \(P\) is already discharged on the present line and should no longer be listed as a primary open problem;
  \item \textbf{absorbed}: \(P\) is no longer independent, because its only remaining unresolved content is already contained in another residual core;
  \item \textbf{residual core}: \(P\) still carries genuinely new proof burden that is not yet discharged by the present line.
\end{enumerate}
\end{definition}

\begin{proposition}[First synchronized placement of the inherited OP queue]
On the present compressed computational line, the inherited named OP queue may be synchronized as follows.
\begin{enumerate}[leftmargin=2em]
  \item \textbf{Closed:} OP~3 is closed as an independent front by the shell-certificate/residual-interface line already exported in the present text, and OP~5 Step~(A) is closed by the source-return, branchwise, and routed carrier-first closure package.
  \item \textbf{Absorbed:} OP~2 is absorbed into OP~1, because the remaining frame-offset/continuum-correction burden no longer survives as an independent Dirac-side front once the present line is synchronized; OP~4 is absorbed into OP~1, because its only still-open quantitative content is precisely the same continuum-limit correction \(\delta_{\mathrm{CL}}\).
  \item \textbf{Residual core:} the live queue is reduced to OP~1 and OP~5 Step~(B). OP~1 remains the continuum-limit/total-scaling residual centered on \(\delta_{\mathrm{CL}}\). OP~5 Step~(B) remains the exact \(\Theta\)--\(R\) balance problem on the routed photon-frame side.
\end{enumerate}
\end{proposition}

\begin{proof}
The earlier OP~3 burden has already been converted into a shell-certified interface theorem together with reusable residual/shell reading rules, so it no longer behaves as an unreduced primary front. Likewise, the source-return closure theorem, its branchwise refinement, and the routed carrier-first OP5 interface already discharge the source/return and admissible-routing part of OP~5, leaving only the exact total balance problem as live burden. Earlier in the text, OP~4 is explicitly reduced to the same continuum-limit correction \(\delta_{\mathrm{CL}}\) as OP~1, and the remaining OP~2 frame-offset burden has already collapsed onto that same continuum-limit core rather than surviving as an independent front. Hence the present synchronized ledger has only two genuine residual cores left: OP~1 and OP~5 Step~(B).
\end{proof}

\begin{corollary}[Plain-language OP synchronization]
The inherited OP queue is no longer a five-front queue. On the current line, OP~3 is closed, OP~2 and OP~4 are absorbed into OP~1, OP~5 survives only through Step~(B), and the live queue is therefore a two-core queue: OP~1 and OP~5 Step~(B).
\end{corollary}

\begin{remark}[Why this synchronization matters]
This synchronization keeps the continuation line honest in both directions. It prevents already discharged or absorbed fronts from being carried forever as if nothing had changed, but it also refuses to pretend that the two genuine residual cores are already gone. The next real version target is therefore sharply focused rather than historically bloated.
\end{remark}

\subsection*{19AH$^{(10)}$. Two-core OP-closing roadmap before the later freeze deferral}
\noindent\textbf{Reader cue.} Read the next roadmap as an earlier compression snapshot of the line. Its purpose is to orient burden genealogy by recording how the queue had once been narrowed to two cores before the later freeze deferral; it does not override the present post-freeze priority order.

After the synchronized ledger from 19AH$^{(9)}$, the line first compressed the next-version target to a narrow two-core closure program. At that stage, the intended next version was to try to close the remaining queue without reopening the green observer-side ledger except where a genuine yellow or red trigger is produced by the re-promotion rule. The purpose is therefore not to add another diagnostic family, but to turn the present compression into two direct closure attempts: one on the continuum-limit residual and one on the exact routed photon-frame balance.

\begin{definition}[Residual-core closing roadmap]
A \emph{residual-core closing roadmap} in this earlier two-core reading is an ordered list of closure tasks that leaves all green observer-side knot families parked by default and reopens them only if a later step actually produces a re-promotion trigger.
\end{definition}

\begin{proposition}[First two-core OP-closing roadmap before the later freeze deferral]
In that earlier two-core reading, the next version would have proceeded through the following tasks.
\begin{enumerate}[leftmargin=2em]
  \item \textbf{OP~1 isolation.} Reduce OP~1 to one final explicit continuum-limit residual identity for \(\delta_{\mathrm{CL}}\), so that no auxiliary historical formulation remains in the live queue.
  \item \textbf{OP~1 closure attempt.} Attack the vanishing or exact identification of that residual directly from the retained total-channel data and the lattice-to-continuum transport structure already built into the present line.
  \item \textbf{OP~5(B) isolation.} Rewrite OP~5 Step~(B) as one exact routed balance statement for the remaining \(\Theta\)--\(R\) defect term, with all source-return, branchwise, and carrier-first routing pieces treated as fixed infrastructure rather than reopened.
  \item \textbf{OP~5(B) closure attempt.} Prove or refute the vanishing of the remaining routed balance defect without importing a second primitive total-channel carrier.
  \item \textbf{Final synchronization pass.} If the two closure attempts succeed, update the synchronized OP ledger so that no live residual core remains; if one attempt fails, record precisely which single residual core survives and keep the green observer-side queue parked.
\end{enumerate}
\end{proposition}

\begin{corollary}[Roadmap discipline in the earlier two-core reading]
In that earlier two-core reading, shell/menu/window, phase-correction, frame-flicker, and other green observer-side families are not first targets. They may be revisited only if closing OP~1 or OP~5 Step~(B) actually forces a yellow or red re-promotion.
\end{corollary}

\begin{remark}[Why the roadmap is intentionally narrow]
The synchronized ledger shows that the continuation line is now past the stage where every historical knot needs to be touched again. The next mathematically serious step in that earlier reading was to spend the new leverage exactly where the residual burden still lives: on the continuum-limit source residual of OP~1 and on the exact routed balance residual of OP~5 Step~(B). A narrower roadmap is therefore not a loss of ambition, but a sign that the line has finally compressed enough to attack the remaining core directly.
\end{remark}




\subsection*{19AH$^{(11)}$. Two-core residual dictionary opening the v\docversion\ closure line}
After the synchronized OP ledger from 19AH$^{(9)}$ and the next-version roadmap from 19AH$^{(10)}$, the correct opening move for v\docversion\ is to stop carrying OP~1 and OP~5 Step~(B) as diffuse historical labels and to rewrite them as two explicit residual objects. Every later closure argument can then be judged by one simple rule: it either changes one of these residuals directly, or it is auxiliary infrastructure only.

\begin{definition}[Continuum-limit residual for OP~1]
\label{def:delta-cl-residual}
Define the \emph{continuum-limit residual}
\[
  \Delta_{\mathrm{CL}}
  :=
  \log \delta_{\mathrm{CL}} + \frac{\gamma_1}{4\pi\gamma_L^5}.
\]
Equivalently, \(\Delta_{\mathrm{CL}}=0\) if and only if the corrected continuum-limit formula from Remark~\ref{conj:cl-correction},
\[
  \delta_{\mathrm{CL}} = \exp\!\left(-\frac{\gamma_1}{4\pi\gamma_L^5}\right),
\]
holds exactly.
\end{definition}

\begin{definition}[Routed photon-frame balance defect for OP~5 Step~(B)]
\label{def:delta-theta-r-residual}
Let \(T:=\gamma_L^{-3/2}\) and let \(\mathcal T(T)\) be the photon-frame Theta sum from Definition~\ref{def:theta-sum-B}. Define the \emph{routed photon-frame balance defect}
\[
  \Delta_{\Theta R}(T)
  :=
  \mathcal T(T)
  - \left(\frac{3T}{2} - \log T - \frac{T^2}{24} - \gamma_E - \bsym^2 T^2\right).
\]
Equivalently, \(\Delta_{\Theta R}(T)=0\) if and only if the exact Theta-sum balance conjecture of Proposition~\ref{conj:theta-sum-B} holds at the routed photon-frame cutoff.
\end{definition}

\begin{proposition}[Two-core reduction of the live OP queue]
\label{prop:two-core-residual-dictionary}
On the synchronized ledger of 19AH$^{(9)}$, the live OP queue is reduced exactly as follows.
\begin{enumerate}[leftmargin=2em]
  \item OP~1 is equivalent to the vanishing or exact first-principles identification of the single scalar residual \(\Delta_{\mathrm{CL}}\).
  \item OP~5 Step~(B) is equivalent to the vanishing of the single routed balance defect \(\Delta_{\Theta R}(T)\).
  \item Any inherited continuation task that leaves both \(\Delta_{\mathrm{CL}}\) and \(\Delta_{\Theta R}(T)\) unchanged remains auxiliary unless it triggers a yellow or red re-promotion in the sense of 19AH$^{(7)}$--19AH$^{(8)}$.
\end{enumerate}
\end{proposition}

\begin{proof}
The OP~1 side was already compressed in theorem form by Theorem~\ref{thm:op1op4-unified}, which reduces the open quantitative content of OP~1 and OP~4 to the single continuum-limit correction \(\delta_{\mathrm{CL}}\). Writing that remaining burden as \(\Delta_{\mathrm{CL}}\) only turns the same statement into one scalar residual identity. On the OP~5 side, Step~(A) is already closed by Proposition~\ref{prop:theta-remainder-trivial}, while Step~(B) is explicitly reduced in Proposition~\ref{conj:theta-sum-B} to one exact routed photon-frame Theta-sum balance statement. Writing that remaining burden as \(\Delta_{\Theta R}(T)\) therefore changes no mathematical content; it only compresses the live queue into one explicit defect variable. The third clause is then the direct reading of the synchronized ledger 19AH$^{(9)}$, the roadmap 19AH$^{(10)}$, and the re-promotion discipline 19AH$^{(7)}$--19AH$^{(8)}$.
\end{proof}

\begin{remark}[Why the new version should start here]
The point of this reduction is not rhetorical neatness. Once the queue is written as \(\Delta_{\mathrm{CL}}\) and \(\Delta_{\Theta R}(T)\), every later step can be audited immediately: either it shrinks one of these two residuals, identifies one of them exactly, or it belongs to infrastructure/support rather than to the remaining core closure burden. This prevents the next version from drifting back into a diffuse many-front reading of the inherited archive.
\end{remark}


\subsection*{19AH$^{(12)}$. First OP~1 slope certificate at the continuum-limit residual}
After the two-core reduction from 19AH$^{(11)}$, the next honest frontal attack on OP~1 is to compress the residual once more: not into a new conjectural formula, but into one normalized attenuation index. This separates the ``what has to be derived'' question from the ``how much surrounding continuum-limit story one tells'' question. The live OP~1 burden is then no longer an undifferentiated derivation of \(\delta_{\mathrm{CL}}\), but the identification of one single scalar exponent/slope quantity.

\begin{definition}[Normalized continuum-limit attenuation index]
\label{def:xi-cl}
Define the \emph{normalized continuum-limit attenuation index}
\[
  \Xi_{\mathrm{CL}}
  :=
  -\frac{4\pi\gamma_L^5}{\gamma_1}\log \delta_{\mathrm{CL}}.
\]
Using Theorem~\ref{thm:op1op4-unified}, equivalently
\[
  \Xi_{\mathrm{CL}}
  =
  -\frac{4\pi\gamma_L^5}{\gamma_1}\log\!\bigl(\sqrt{3}\,\beta_{\mathrm{sym}}\bigr).
\]
Thus \(\Xi_{\mathrm{CL}}\) measures exactly how far the still-open continuum-limit correction deviates from the candidate exponent in Remark~\ref{conj:cl-correction}, but in normalized unit form.
\end{definition}

\begin{proposition}[Equivalent slope certificate for OP~1]
\label{prop:op1-slope-certificate}
The following statements are equivalent.
\begin{enumerate}[leftmargin=2em]
  \item The OP~1 continuum-limit residual vanishes, i.e. \(\Delta_{\mathrm{CL}}=0\).
  \item The normalized attenuation index satisfies \(\Xi_{\mathrm{CL}}=1\).
  \item Whenever one can write the remaining continuum-limit correction at the first live scale in exponential form
  \[
    \delta_{\mathrm{CL}} = e^{-c_{\mathrm{CL}}\gamma_1},
  \]
  for some dimensionless attenuation slope \(c_{\mathrm{CL}}\), that slope is exactly
  \[
    c_{\mathrm{CL}} = \frac{1}{4\pi\gamma_L^5}.
  \]
\end{enumerate}
\end{proposition}

\begin{proof}
By Definition~\ref{def:delta-cl-residual},
\[
  \Delta_{\mathrm{CL}}=0
  \quad\Longleftrightarrow\quad
  \log \delta_{\mathrm{CL}} = -\frac{\gamma_1}{4\pi\gamma_L^5}.
\]
Multiplying by \(-4\pi\gamma_L^5/\gamma_1\) gives
\[
  \Delta_{\mathrm{CL}}=0
  \quad\Longleftrightarrow\quad
  -\frac{4\pi\gamma_L^5}{\gamma_1}\log \delta_{\mathrm{CL}} = 1,
\]
which is exactly \(\Xi_{\mathrm{CL}}=1\), proving the equivalence of (1) and (2). If in addition \(\delta_{\mathrm{CL}} = e^{-c_{\mathrm{CL}}\gamma_1}\), then
\[
  \log \delta_{\mathrm{CL}} = -c_{\mathrm{CL}}\gamma_1,
\]
so the same identity becomes
\[
  c_{\mathrm{CL}}\gamma_1 = \frac{\gamma_1}{4\pi\gamma_L^5}
  \quad\Longleftrightarrow\quad
  c_{\mathrm{CL}} = \frac{1}{4\pi\gamma_L^5}.
\]
Conversely, this slope value immediately reconstructs the vanishing residual formula from Definition~\ref{def:delta-cl-residual}. Hence all three statements are equivalent.
\end{proof}

\begin{corollary}[What the r02 OP~1 attack reduces to]
\label{cor:op1-r02-target}
On the present compressed line, a successful OP~1 closure attempt no longer needs to present itself as a broad continuum-limit narrative. It is enough to derive, certify, or uniquely force the single normalized attenuation statement
\[
  \Xi_{\mathrm{CL}}=1.
\]
Equivalently, under any admissible exponential first-scale attenuation ansatz, OP~1 is closed once the slope \(c_{\mathrm{CL}}\) is shown to equal \(1/(4\pi\gamma_L^5)\).
\end{corollary}

\begin{remark}[Why this is a real reduction and not just notation]
This step still does \emph{not} prove OP~1. What it does remove is one layer of ambiguity about what the next version must actually derive. The live target is no longer a vague \emph{prime-lattice / continuum-limit correction story}; it is the single normalized exponent encoded by \(\Xi_{\mathrm{CL}}\). Any later lemma about lattice counting, prime spacing, frame conversion, or first-zero transport is now auditable by one question only: does it fix \(\Xi_{\mathrm{CL}}\), leave it unchanged, or force it away from~1? In that sense r02 turns OP~1 into the narrowest honest one-scalar closure gate currently available.
\end{remark}


\subsection*{19AH$^{(13)}$. First frame-transport split certificate for OP~1 in photon-first-zero normalization}
After the slope certificate from 19AH$^{(12)}$, the next useful compression is not to invent a new exponent but to split the already isolated OP~1 load into its frame-conversion part and its photon-side transport part. Remark~\ref{conj:cl-correction} already records that the candidate residual formula may be written in two equivalent ways,
\[
  \delta_{\mathrm{CL}}
  = \exp\!\left(-\frac{\gamma_1}{4\pi\gamma_L^5}\right)
  = \exp\!\left(-\frac{\gamma_1^{(\mathrm{photon})}}{4\pi\gamma_L^4}\right),
\]
with \(\gamma_1^{(\mathrm{photon})}=\gamma_1/\gamma_L\). The fifth power of \(\gamma_L\) in the matter-frame formula is therefore already visibly split into one frame-change factor and four photon-side transport factors. This means that OP~1 can be attacked in a sharper way: not as a mysterious global \(\gamma_L^5\)-law, but as a photon-first-zero transport law with four residual shell factors after the frame conversion has been peeled off.

\begin{definition}[Photon-normalized continuum-transport index]
\label{def:psi-cl}
Define the \emph{photon-normalized continuum-transport index}
\[
  \Psi_{\mathrm{CL}}
  :=
  -\frac{4\pi\gamma_L^4}{\gamma_1^{(\mathrm{photon})}}\log \delta_{\mathrm{CL}}.
\]
Equivalently, since \(\gamma_1^{(\mathrm{photon})}=\gamma_1/\gamma_L\),
\[
  \Psi_{\mathrm{CL}}
  =
  -\frac{4\pi\gamma_L^5}{\gamma_1}\log \delta_{\mathrm{CL}}
  = \Xi_{\mathrm{CL}}.
\]
Thus \(\Psi_{\mathrm{CL}}\) carries the same open numerical burden as \(\Xi_{\mathrm{CL}}\), but written after the matter-to-photon frame passage has been made explicit.
\end{definition}

\begin{proposition}[Frame-transport split certificate for OP~1]
\label{prop:op1-frame-transport-split}
The following statements are equivalent.
\begin{enumerate}[leftmargin=2em]
  \item The OP~1 residual vanishes, i.e. \(\Delta_{\mathrm{CL}}=0\).
  \item The slope certificate of 19AH$^{(12)}$ holds, i.e. \(\Xi_{\mathrm{CL}}=1\).
  \item The photon-normalized transport index satisfies \(\Psi_{\mathrm{CL}}=1\).
  \item The continuum-limit correction has the photon-first-zero transport form
  \[
    \delta_{\mathrm{CL}}
    = \exp\!\left(-\frac{\gamma_1^{(\mathrm{photon})}}{4\pi\gamma_L^4}\right).
  \]
\end{enumerate}
In particular, OP~1 is equivalent to proving that the entire residual attenuation beyond the frame change from \(\gamma_1\) to \(\gamma_1^{(\mathrm{photon})}\) is exactly one unit of photon-normalized four-layer transport load.
\end{proposition}

\begin{proof}
The equivalence of (1) and (2) is Proposition~\ref{prop:op1-slope-certificate}. By Definition~\ref{def:psi-cl} and the identity \(\gamma_1^{(\mathrm{photon})}=\gamma_1/\gamma_L\), one has
\[
  \Psi_{\mathrm{CL}}
  =
  -\frac{4\pi\gamma_L^4}{\gamma_1^{(\mathrm{photon})}}\log \delta_{\mathrm{CL}}
  =
  -\frac{4\pi\gamma_L^5}{\gamma_1}\log \delta_{\mathrm{CL}}
  = \Xi_{\mathrm{CL}},
\]
so (2) and (3) are equivalent. Solving \(\Psi_{\mathrm{CL}}=1\) for \(\delta_{\mathrm{CL}}\) gives
\[
  \log \delta_{\mathrm{CL}} = -\frac{\gamma_1^{(\mathrm{photon})}}{4\pi\gamma_L^4},
\]
and exponentiating yields (4). Conversely, (4) inserted into Definition~\ref{def:psi-cl} gives \(\Psi_{\mathrm{CL}}=1\). Hence all four statements are equivalent.
\end{proof}

\begin{corollary}[What the r03 OP~1 attack really changes]
\label{cor:op1-r03-target}
The remaining OP~1 burden may now be read in a more structured way. One no longer has to ``explain a \(\gamma_L^5\) exponent all at once''. It is enough to justify the following split picture:
\begin{enumerate}[leftmargin=2em]
  \item one power of \(\gamma_L\) comes from the explicit passage from the matter-frame first zero \(\gamma_1\) to the photon-frame first zero \(\gamma_1^{(\mathrm{photon})}\);
  \item the remaining four powers form the actual photon-side transport attenuation burden;
  \item that photon-side burden has normalized weight exactly one.
\end{enumerate}
Therefore any future counting, lattice, shell, or first-zero transport lemma may now be audited by the sharper question: does it explain the unit value of \(\Psi_{\mathrm{CL}}\), rather than merely restating the matter-frame exponent?
\end{corollary}

\begin{remark}[Why this is progress without pretending closure]
This still does not prove OP~1. The real gain is conceptual localization. In r02 the target was the one-scalar index \(\Xi_{\mathrm{CL}}\). In r03 the same scalar is re-read in a frame-split form that matches the standing photon/matter distinction of the Companion. That removes one source of merge-style ambiguity: the fifth power of \(\gamma_L\) no longer needs to be read as an indivisible miracle. The present line may now search specifically for a \emph{four-layer photon transport law after one explicit frame conversion}. That is narrower, structurally cleaner, and closer to the actual lattice/transport heuristics already used elsewhere in the document.
\end{remark}

\subsection*{19AH$^{(14)}$. First finite four-shell balance certificate for OP~1 after the frame split}
After the frame-transport split of 19AH$^{(13)}$, the next honest compression is to stop carrying the photon-side transport burden as one undifferentiated four-power block and to write it as a finite shell balance. This does not yet explain \emph{why} the four-layer photon transport should have normalized weight one, but it does isolate the exact finite statement that any such explanation would have to prove. In particular, OP~1 can now be re-read as a question about the total normalized load carried by four post-frame transport layers, rather than as a raw continuum-limit estimate.

\begin{definition}[Finite four-shell transport loads for the OP~1 residual]
\label{def:op1-four-shell-loads}
Assume the photon-side continuum-limit attenuation admits a four-factor decomposition
\[
  \delta_{\mathrm{CL}} = \prod_{j=1}^4 \tau_j,
\]
with positive shell factors \(\tau_j>0\). Define the corresponding \emph{normalized shell loads}
\[
  \lambda_j := -\frac{4\pi\gamma_L^4}{\gamma_1^{(\mathrm{photon})}}\log \tau_j
  \qquad (j=1,2,3,4).
\]
Whenever the above product formula holds, the photon-normalized transport index of Definition~\ref{def:psi-cl} satisfies
\[
  \Psi_{\mathrm{CL}} = \sum_{j=1}^4 \lambda_j.
\]
\end{definition}

\begin{proposition}[Finite four-shell closure certificate for OP~1]
\label{prop:op1-four-shell-certificate}
Suppose \(\delta_{\mathrm{CL}}=\prod_{j=1}^4 \tau_j\) is a positive four-shell decomposition and let \(\lambda_j\) be the normalized shell loads from Definition~\ref{def:op1-four-shell-loads}. Then the following are equivalent.
\begin{enumerate}[leftmargin=2em]
  \item The OP~1 residual vanishes, i.e. \(\Delta_{\mathrm{CL}}=0\).
  \item The photon-normalized transport index satisfies \(\Psi_{\mathrm{CL}}=1\).
  \item The four normalized shell loads satisfy the finite balance law
  \[
    \sum_{j=1}^4 \lambda_j = 1.
  \]
\end{enumerate}
In particular, after the frame split from 19AH$^{(13)}$, OP~1 is equivalent to proving that the post-frame four-shell transport carries total normalized load exactly one.
\end{proposition}

\begin{proof}
By Proposition~\ref{prop:op1-frame-transport-split}, statement (1) is equivalent to \(\Psi_{\mathrm{CL}}=1\), so (1) and (2) are equivalent. If \(\delta_{\mathrm{CL}}=\prod_{j=1}^4 \tau_j\), then
\[
  \log \delta_{\mathrm{CL}} = \sum_{j=1}^4 \log \tau_j.
\]
Multiplying by \(-4\pi\gamma_L^4/\gamma_1^{(\mathrm{photon})}\) gives
\[
  \Psi_{\mathrm{CL}}
  = -\frac{4\pi\gamma_L^4}{\gamma_1^{(\mathrm{photon})}}\log \delta_{\mathrm{CL}}
  = \sum_{j=1}^4 \left(-\frac{4\pi\gamma_L^4}{\gamma_1^{(\mathrm{photon})}}\log \tau_j\right)
  = \sum_{j=1}^4 \lambda_j,
\]
which is exactly statement (3). Hence (2) and (3) are equivalent, and all three assertions are equivalent.
\end{proof}

\begin{corollary}[Equal-shell partition is sufficient but not necessary]
\label{cor:op1-quarter-shells}
Under the hypotheses of Proposition~\ref{prop:op1-four-shell-certificate}, any of the following is sufficient to close OP~1.
\begin{enumerate}[leftmargin=2em]
  \item \(\lambda_j=\tfrac14\) for each \(j=1,2,3,4\);
  \item more generally, any asymmetric four-shell profile with
  \[
    \lambda_1+\lambda_2+\lambda_3+\lambda_4 = 1.
  \]
\end{enumerate}
Thus the present line does \emph{not} require equal quarter-load shells; equal partition is merely the simplest admissible closure pattern.
\end{corollary}



\begin{remark}[Why r04 is a genuine narrowing of OP~1]
This still does not prove OP~1, and it should not be oversold as such. The genuine gain is that OP~1 is now finite twice over. First, r02 reduced it to one scalar index; second, r03 split that index into one frame factor plus a four-layer photon burden; now r04 turns the surviving photon burden into a finite additive balance law. The live question is no longer \emph{``what is the whole continuum-limit correction story?''} but rather: \emph{which four post-frame shell loads are present, and why do they sum to one?} Any future lattice, shell, counting, or transport lemma can therefore be audited even more sharply: does it produce one of the \(\lambda_j\), constrain their sum, or leave the finite balance untouched?
\end{remark}

\subsection*{19AH$^{(15)}$. First two-pass half-burden certificate for the OP~1 four-shell balance}
The finite four-shell balance of 19AH$^{(14)}$ still leaves open where the four loads \(\lambda_j\) should come from. The first structurally natural candidate is a symmetrized two-pass reading of the post-frame photon transport: two shell loads belong to the outward/forward pass and two belong to the return/backward pass. This does not yet determine the individual shell factors, but it reduces the origin question from four unrelated loads to two pass totals. That is already enough to expose the first genuinely bidirectional closure pattern for OP~1.

\begin{definition}[Two-pass shell-pair loads for the OP~1 residual]
\label{def:op1-two-pass-loads}
Under the hypotheses of Definition~\ref{def:op1-four-shell-loads}, define the forward- and backward-pass totals
\[
  \Lambda_{\rightarrow} := \lambda_1+\lambda_2,
  \qquad
  \Lambda_{\leftarrow} := \lambda_3+\lambda_4.
\]
We call the four-shell transport profile \emph{two-pass grouped} if the post-frame shell loads are read as two forward-side loads and two backward-side loads in this way. We call it \emph{pass-symmetric} if in addition
\[
  \Lambda_{\rightarrow}=\Lambda_{\leftarrow}.
\]
\end{definition}

\begin{proposition}[Two-pass half-burden certificate for OP~1]
\label{prop:op1-two-pass-half-burden}
Suppose the OP~1 residual admits a two-pass grouped four-shell decomposition in the sense of Definition~\ref{def:op1-two-pass-loads}. Then the following are equivalent.
\begin{enumerate}[leftmargin=2em]
  \item The OP~1 residual vanishes, i.e. \(\Delta_{\mathrm{CL}}=0\).
  \item The photon-side transport index satisfies \(\Psi_{\mathrm{CL}}=1\).
  \item The two pass totals satisfy
  \[
    \Lambda_{\rightarrow}+\Lambda_{\leftarrow}=1.
  \]
\end{enumerate}
If the profile is moreover pass-symmetric, then these are also equivalent to
\[
  \Lambda_{\rightarrow}=\Lambda_{\leftarrow}=\frac12.
\]
Thus, under pass symmetry, OP~1 is equivalent to saying that the outward and return photon-side passes each carry normalized half-burden.
\end{proposition}

\begin{proof}
By Proposition~\ref{prop:op1-four-shell-certificate}, \(\Delta_{\mathrm{CL}}=0\) is equivalent to
\[
  \lambda_1+\lambda_2+\lambda_3+\lambda_4=1.
\]
Using Definition~\ref{def:op1-two-pass-loads}, this is exactly
\[
  \Lambda_{\rightarrow}+\Lambda_{\leftarrow}=1,
\]
so statements (1)--(3) are equivalent. If, in addition, the profile is pass-symmetric, then \(\Lambda_{\rightarrow}=\Lambda_{\leftarrow}\); together with their sum being one, this gives
\[
  \Lambda_{\rightarrow}=\Lambda_{\leftarrow}=\tfrac12.
\]
Conversely, if both pass totals equal \(\tfrac12\), then their sum is one, and OP~1 closes by the already established equivalences.
\end{proof}

\begin{corollary}[What r05 really reduces]
\label{cor:op1-two-pass-reduction}
Under a pass-symmetric two-pass reading, OP~1 no longer needs a direct identification of all four individual shell loads. It is enough to show that the forward pass contributes normalized burden \(\tfrac12\) and the backward pass contributes normalized burden \(\tfrac12\). In particular, any future lattice, first-zero, shell, or transport lemma that yields pass symmetry together with one half-burden identity closes OP~1 without resolving the internal split inside each pass.
\end{corollary}

\begin{remark}[Why this is the first honest origin candidate for the four \texorpdfstring{$\lambda_j$}{lambdaj}]
This is not yet a first-principles derivation, and it should not be advertised as one. The gain is narrower and more honest. The four shell loads now have a concrete structural reading as \emph{two forward-side loads plus two backward-side loads}. That matches the standing bidirectional transport intuition already present elsewhere in the project and turns OP~1 from a raw four-term balance into a two-pass half-burden question. The next attack can therefore be sharper: not ``guess the four \(\lambda_j\) independently,'' but derive why the post-frame photon transport should be pass-symmetric and why each pass should carry normalized load \(\tfrac12\).
\end{remark}


\subsection*{19AH$^{(16)}$. First reciprocal pass-symmetry certificate for the OP~1 two-pass balance}
After the two-pass half-burden certificate of 19AH$^{(15)}$, the next narrow question is no longer how to guess the four shell loads independently, but why the forward and backward pass totals should agree in the first place. The weakest structurally natural answer is a reciprocity principle: each forward-side shell load has a backward-side partner carrying the same normalized attenuation burden. This does not yet derive the common value of the pairs, but it is enough to force pass symmetry and therefore collapse OP~1 to a one-pass half-balance statement.

\begin{definition}[Reciprocal shell-pairing certificate for the OP~1 transport]
\label{def:op1-reciprocal-shell-pairing}
Under the hypotheses of Definition~\ref{def:op1-two-pass-loads}, we say that the two-pass grouped four-shell profile satisfies the \emph{reciprocal shell-pairing certificate} if the forward/backward shell labels are paired by
\[
  \rho(1)=3,\qquad \rho(2)=4,
\]
and the paired normalized loads agree,
\[
  \lambda_{\rho(j)}=\lambda_j\qquad (j=1,2).
\]
Equivalently,
\[
  \lambda_1=\lambda_3,
  \qquad
  \lambda_2=\lambda_4.
\]
We call this \emph{pairwise reciprocal} because each outward shell burden is mirrored by one return-shell burden of the same normalized weight.
\end{definition}

\begin{proposition}[Reciprocal pairing forces pass symmetry for OP~1]
\label{prop:op1-reciprocal-pass-symmetry}
Suppose the OP~1 residual admits a two-pass grouped four-shell decomposition in the sense of Definition~\ref{def:op1-two-pass-loads}. If the reciprocal shell-pairing certificate of Definition~\ref{def:op1-reciprocal-shell-pairing} holds, then the profile is automatically pass-symmetric:
\[
  \Lambda_{\rightarrow}=\Lambda_{\leftarrow}.
\]
Consequently,
\[
  \Delta_{\mathrm{CL}}=0
  \iff
  \Psi_{\mathrm{CL}}=1
  \iff
  \Lambda_{\rightarrow}=\Lambda_{\leftarrow}=\frac12.
\]
Equivalently, under reciprocal shell pairing, OP~1 is reduced to the single one-pass half-burden identity
\[
  \lambda_1+\lambda_2=\frac12,
\]
or, identically, to the return-pass statement
\[
  \lambda_3+\lambda_4=\frac12.
\]
\end{proposition}

\begin{proof}
If reciprocal shell pairing holds, then
\[
  \lambda_1=\lambda_3,
  \qquad
  \lambda_2=\lambda_4.
\]
Therefore
\[
  \Lambda_{\rightarrow}=\lambda_1+\lambda_2
  =\lambda_3+\lambda_4=\Lambda_{\leftarrow},
\]
so the profile is pass-symmetric. The remaining equivalences are then exactly the pass-symmetric conclusion of Proposition~\ref{prop:op1-two-pass-half-burden}. In particular, once reciprocal pairing is known, OP~1 closes precisely when one of the two pass totals equals \(\tfrac12\), because the other is then forced to equal the same value.
\end{proof}

\begin{corollary}[The live OP~1 burden after r06]
\label{cor:op1-live-burden-after-r06}
If a future shell, lattice, counting, or first-zero transport lemma yields reciprocal shell pairing, then the separate proof obligation for pass symmetry disappears. The live OP~1 burden is then only the one-pass half-balance statement
\[
  \lambda_1+\lambda_2=\frac12,
\]
with the backward pass following automatically by reciprocity.
\end{corollary}

\begin{remark}[Why this is a real narrowing without pretending OP~1 is closed]
This is still not a first-principles derivation of the continuum-limit correction, and it should not be sold as such. The gain is more disciplined. After r05, OP~1 asked for two ingredients: pass symmetry and half-burden. After r06, any concrete reciprocity mechanism collapses the first ingredient into a taut consequence of pairwise equal shell loads. The next attack can therefore become even sharper: not ``why should the two passes be equal in aggregate?'' but ``what concrete structural mechanism enforces reciprocal shell pairing, and why does one pass carry normalized load \(\tfrac12\)?''
\end{remark}


\subsection*{19AH$^{(17)}$. First same-corridor reversal certificate for reciprocal shell pairing in OP~1}
After the reciprocal pass-symmetry certificate of 19AH$^{(16)}$, the next honest question is what concrete structural mechanism could force the pairwise equalities
\[
  \lambda_1=\lambda_3,
  \qquad
  \lambda_2=\lambda_4.
\]
The narrowest bidirectional candidate is the standing two-pass reading itself: the post-frame photon transport uses the \emph{same shell corridor} once in the forward direction and once in the reverse direction. If the normalized attenuation burden is attached to the corridor shells rather than to the direction label, then reciprocal shell pairing follows automatically. This still does not determine the common shell values, but it supplies the first concrete origin mechanism for the r06 reciprocity pattern.

\begin{definition}[Same-corridor reversal certificate for the OP~1 shell transport]
\label{def:op1-same-corridor-reversal}
Assume the grouped four-shell decomposition of Definition~\ref{def:op1-two-pass-loads}. We say that the profile satisfies the \emph{same-corridor reversal certificate} if there exist two nonnegative corridor loads
\[
  \mu_1,\mu_2\ge 0
\]
and a single two-shell corridor such that the forward pass reads the corridor in order
\[
  (\lambda_1,\lambda_2)=(\mu_1,\mu_2),
\]
while the backward pass reads the same corridor in reverse order but with the same shell-attached loads,
\[
  (\lambda_3,\lambda_4)=(\mu_1,\mu_2).
\]
Equivalently, the direction of traversal changes, but the shell-burden assignment is carried by the corridor itself and therefore survives reversal unchanged.
\end{definition}

\begin{proposition}[Same-corridor reversal implies reciprocal shell pairing]
\label{prop:op1-same-corridor-implies-reciprocity}
If the OP~1 grouped four-shell profile satisfies the same-corridor reversal certificate of Definition~\ref{def:op1-same-corridor-reversal}, then the reciprocal shell-pairing certificate of Definition~\ref{def:op1-reciprocal-shell-pairing} holds. Hence
\[
  \lambda_1=\lambda_3=\mu_1,
  \qquad
  \lambda_2=\lambda_4=\mu_2,
\]
and therefore
\[
  \Lambda_{\rightarrow}=\Lambda_{\leftarrow}=\mu_1+\mu_2.
\]
Consequently,
\[
  \Delta_{\mathrm{CL}}=0
  \iff
  \Psi_{\mathrm{CL}}=1
  \iff
  \mu_1+\mu_2=\frac12.
\]
In particular, under same-corridor reversal the full OP~1 burden is reduced to one two-shell half-balance on a single direction-free corridor.
\end{proposition}

\begin{proof}
By definition of same-corridor reversal we have
\[
  \lambda_1=\mu_1=\lambda_3,
  \qquad
  \lambda_2=\mu_2=\lambda_4.
\]
Thus the reciprocal shell-pairing equalities of Definition~\ref{def:op1-reciprocal-shell-pairing} hold automatically. Proposition~\ref{prop:op1-reciprocal-pass-symmetry} then yields pass symmetry, and the pass totals are
\[
  \Lambda_{\rightarrow}=\lambda_1+\lambda_2=\mu_1+\mu_2,
  \qquad
  \Lambda_{\leftarrow}=\lambda_3+\lambda_4=\mu_1+\mu_2.
\]
Since OP~1 is equivalent to \(\Lambda_{\rightarrow}=\Lambda_{\leftarrow}=\tfrac12\) under reciprocal pairing, the final equivalence follows.
\end{proof}

\begin{corollary}[The live OP~1 burden after r07]
\label{cor:op1-live-burden-after-r07}
If a future prime-lattice, shell, or first-zero transport lemma identifies the post-frame photon transport with one reversal-invariant two-shell corridor, then OP~1 no longer asks for a separate reciprocity proof and no longer asks for two independent pass totals. The live burden is only the scalar two-shell corridor identity
\[
  \mu_1+\mu_2=\frac12.
\]
\end{corollary}

\begin{remark}[Why this is the first genuinely concrete reciprocity mechanism]
This is still not the closure of OP~1. The unresolved content has merely been pushed to an even narrower place: the one-corridor half-balance \(\mu_1+\mu_2=\tfrac12\). But unlike the earlier certificates, r07 now gives a concrete structural reason for the reciprocity equalities themselves. The forward and backward passes are no longer treated as two independently weighted routes; they are two traversals of the same post-frame corridor. That is the first point in the present line where the project's standing bidirectional transport picture becomes an explicit OP~1 origin mechanism rather than only a heuristic slogan.
\end{remark}



\subsection*{19AH$^{(18)}$. First unit-loop half-burden certificate for the one-corridor OP~1 transport}
After the same-corridor reversal certificate of 19AH$^{(17)}$, the unresolved OP~1 content has been pushed onto one direction-free two-shell corridor with total pass load \(\mu_1+\mu_2\). The next honest question is therefore no longer where reciprocity comes from, but why this one corridor should carry exactly the half-burden \(\tfrac12\). The narrowest bidirectional candidate is the standing two-pass transport picture itself: the post-frame photon corridor is traversed once forward and once backward, and the normalized burden is attached to the full closed two-pass loop rather than to one preferred orientation. If that closed loop is unit-normalized and the corridor loads are reversal-invariant, then each traversal must carry exactly half of the total loop burden.

\begin{definition}[Unit-loop half-burden certificate for the OP~1 one-corridor transport]
\label{def:op1-unit-loop-half-burden}
Assume the same-corridor reversal certificate of Definition~\ref{def:op1-same-corridor-reversal}, so that a single direction-free two-shell corridor with loads \(\mu_1,\mu_2\ge 0\) is traversed once forward and once backward. We say that the profile satisfies the \emph{unit-loop half-burden certificate} if the full closed two-pass corridor loop is normalized to total burden one, i.e.
\[
  (\mu_1+\mu_2)+(\mu_1+\mu_2)=1.
\]
Equivalently, the forward and backward traversals are not assigned independent total weights; they are the two orientations of one unit-normalized closed transport loop.
\end{definition}

\begin{proposition}[Unit-loop normalization forces the OP~1 half-burden]
\label{prop:op1-unit-loop-half-burden}
If the OP~1 transport satisfies the same-corridor reversal certificate of Definition~\ref{def:op1-same-corridor-reversal} and the unit-loop half-burden certificate of Definition~\ref{def:op1-unit-loop-half-burden}, then
\[
  \mu_1+\mu_2=\frac12.
\]
Consequently,
\[
  \Lambda_{\rightarrow}=\Lambda_{\leftarrow}=\frac12,
\qquad
  \Psi_{\mathrm{CL}}=1,
\qquad
  \Delta_{\mathrm{CL}}=0.
\]
In particular, under same-corridor reversal plus unit-loop normalization, OP~1 is closed.
\end{proposition}

\begin{proof}
By Definition~\ref{def:op1-unit-loop-half-burden},
\[
  2(\mu_1+\mu_2)=1.
\]
Hence \(\mu_1+\mu_2=\tfrac12\). Proposition~\ref{prop:op1-same-corridor-implies-reciprocity} then gives
\[
  \Lambda_{\rightarrow}=\Lambda_{\leftarrow}=\mu_1+\mu_2=\frac12.
\]
By Proposition~\ref{prop:op1-same-corridor-implies-reciprocity}, this is equivalent to \(\Psi_{\mathrm{CL}}=1\), and by Proposition~\ref{prop:op1-slope-certificate} it is equivalent to \(\Delta_{\mathrm{CL}}=0\). Thus OP~1 is closed under the stated hypotheses.
\end{proof}

\begin{corollary}[The live OP~1 burden after r08]
\label{cor:op1-live-burden-after-r08}
After r08, the remaining live content of OP~1 is no longer a free four-shell, two-pass, or reciprocity problem. It is the identification of the post-frame photon corridor with a \emph{unit-normalized closed two-pass loop}. Once that normalization is obtained from a future prime-lattice, shell, or first-zero transport lemma, the one-corridor half-balance and therefore OP~1 follow immediately.
\end{corollary}

\begin{remark}[Why r08 is a real advance without pretending the normalization is already proved]
This does not yet prove that the post-frame photon corridor really is unit-normalized. That point remains the live burden. But the nature of the burden has changed sharply. After r07, OP~1 asked for one scalar half-balance on a direction-free corridor. After r08, even that half-balance is no longer independent: it is forced as soon as the corridor is shown to be the correct unit-normalized closed two-pass object. So the next attack does not need to juggle shell weights, pass symmetry, or reciprocity any more. It only needs to justify one normalization principle for one already isolated bidirectional corridor.
\end{remark}


\subsection*{19AH$^{(19)}$. First one-first-zero loop-budget certificate for OP~1 after the unit-loop reduction}
After the unit-loop half-burden certificate of 19AH$^{(18)}$, the remaining burden is no longer to guess a half-load on the corridor. What remains is to explain why the isolated closed two-pass corridor should carry \emph{one} normalized transport unit in the first place. The frame split from 19AH$^{(13)}$ already shows that the surviving OP~1 normalization is measured in units of the photon-normalized first zero \(\gamma_1^{(\mathrm{photon})}\). The sharpest next compression is therefore to re-read the unit-loop hypothesis not as an external normalization convention, but as a multiplicity-one statement: the post-frame corridor is exactly one closed first-zero transport loop, not a fraction of one loop and not a multiple copy of several such loops.

\begin{definition}[One-first-zero loop-budget certificate for the OP~1 corridor]
\label{def:op1-one-first-zero-loop-budget}
Assume the same-corridor reversal certificate of Definition~\ref{def:op1-same-corridor-reversal}, so that the post-frame OP~1 transport is carried by one direction-free two-shell corridor with loads \(\mu_1,\mu_2\). Define its \emph{closed-loop budget index} by
\[
  \Omega_{\mathrm{loop}} := 2(\mu_1+\mu_2).
\]
We say that the profile satisfies the \emph{one-first-zero loop-budget certificate} if the isolated closed two-pass corridor represents exactly one photon-first-zero transport unit, i.e.
\[
  \Omega_{\mathrm{loop}} = 1.
\]
Equivalently, the corridor is read as a multiplicity-one closed loop in the photon-first-zero normalization of 19AH$^{(13)}$.
\end{definition}

\begin{proposition}[One-first-zero loop budget identifies the r08 unit-loop normalization]
\label{prop:op1-one-first-zero-loop-budget}
Assume the same-corridor reversal certificate of Definition~\ref{def:op1-same-corridor-reversal}. Then
\[
  \Omega_{\mathrm{loop}} = \Psi_{\mathrm{CL}}.
\]
Consequently, if the one-first-zero loop-budget certificate of Definition~\ref{def:op1-one-first-zero-loop-budget} holds, then the unit-loop half-burden certificate of Definition~\ref{def:op1-unit-loop-half-burden} holds as well, and therefore
\[
  \Delta_{\mathrm{CL}} = 0.
\]
In particular, under same-corridor reversal, OP~1 is closed as soon as the isolated corridor is shown to be a multiplicity-one photon-first-zero loop.
\end{proposition}

\begin{proof}
By Proposition~\ref{prop:op1-same-corridor-implies-reciprocity},
\[
  \Lambda_{\rightarrow}=\Lambda_{\leftarrow}=\mu_1+\mu_2.
\]
Hence Proposition~\ref{prop:op1-two-pass-half-burden} gives
\[
  \Psi_{\mathrm{CL}} = \Lambda_{\rightarrow}+\Lambda_{\leftarrow}
  = (\mu_1+\mu_2)+(\mu_1+\mu_2)
  = 2(\mu_1+\mu_2)
  = \Omega_{\mathrm{loop}}.
\]
If now \(\Omega_{\mathrm{loop}}=1\), then by definition
\[
  (\mu_1+\mu_2)+(\mu_1+\mu_2)=1,
\]
which is exactly the unit-loop half-burden certificate of Definition~\ref{def:op1-unit-loop-half-burden}. Proposition~\ref{prop:op1-unit-loop-half-burden} then yields \(\Delta_{\mathrm{CL}}=0\).
\end{proof}

\begin{corollary}[The live OP~1 burden after r09]
\label{cor:op1-live-burden-after-r09}
After r09, the remaining live burden of OP~1 is no longer to supply an abstract unit-loop normalization from outside the frame split. It is to prove a \emph{multiplicity-one statement}: the isolated post-frame corridor is exactly one photon-first-zero closed loop. Equivalently, one must rule out both undercounting and overcounting, i.e. show that the corridor does not represent a fractional share of the first-zero budget and does not split into several loop copies carrying total index \(\Omega_{\mathrm{loop}}\neq 1\).
\end{corollary}

\begin{remark}[Why r09 is a real tightening of the OP~1 queue]
This is still not the final closure proof, because the multiplicity-one statement has not yet been derived from the prime-lattice/first-zero transport itself. But r09 removes one remaining piece of arbitrariness. In r08, the decisive hypothesis was a unit-normalized closed loop. After r09, that hypothesis is no longer a free normalization slogan: it is identified with the claim that the already isolated corridor carries exactly one photon-first-zero transport unit. So the next attack on OP~1 is even narrower than before. It should no longer ask for ``why does some loop have weight one?'' in the abstract, but for a concrete first-zero transport reason why the admissible corridor is a single normalized loop rather than a fractional or multiple-count loop.
\end{remark}



\subsection*{19AH$^{(20)}$. First admissible-window uniqueness certificate for the multiplicity-one OP~1 corridor}
After the one-first-zero loop-budget certificate of 19AH$^{(19)}$, the live OP~1 burden is no longer a free normalization question. It is a uniqueness question on the retained admissible corridor: why should the isolated post-frame same-corridor loop represent \emph{exactly one} admissible photon-first-zero transport unit rather than several parallel admissible copies? The sharpest next compression is therefore to move from multiplicity language to admissible-window uniqueness language. On the present line, this is the first point at which the older no-second-carrier discipline and the newer total/imbalance burden split directly feed back into OP~1: extra loop multiplicity can only matter if it either changes the retained total-channel profile or introduces a genuinely new primitive carrier contribution on the same admissible corridor.

\begin{definition}[Admissible-window uniqueness certificate for the OP~1 corridor]
\label{def:op1-admissible-window-uniqueness}
Assume the same-corridor reversal certificate of Definition~\ref{def:op1-same-corridor-reversal}, so that the post-frame OP~1 transport is carried by one direction-free two-shell corridor. We say that the profile satisfies the \emph{admissible-window uniqueness certificate} if, on the intended admissible machine-bearing windows of the present line, that isolated corridor admits exactly one primitive photon-first-zero closed-loop realization compatible with the retained total-channel profile and with the no-second-carrier discipline. Equivalently, every admissible realization of the corridor that preserves the retained total-channel load is identified with the same single primitive first-zero loop.
\end{definition}

\begin{proposition}[Admissible-window uniqueness upgrades multiplicity-one to OP~1 closure]
\label{prop:op1-admissible-window-uniqueness}
Assume the same-corridor reversal certificate of Definition~\ref{def:op1-same-corridor-reversal}. If the admissible-window uniqueness certificate of Definition~\ref{def:op1-admissible-window-uniqueness} holds, then the one-first-zero loop-budget certificate of Definition~\ref{def:op1-one-first-zero-loop-budget} holds. Consequently,
\[
  \Delta_{\mathrm{CL}} = 0.
\]
In particular, after r10 the surviving OP~1 burden may be read as a pure admissible-window uniqueness question for the retained post-frame corridor.
\end{proposition}

\begin{proof}
By Definition~\ref{def:op1-admissible-window-uniqueness}, the isolated post-frame same-corridor loop has exactly one admissible primitive photon-first-zero realization compatible with the retained total-channel profile and the no-second-carrier discipline. But Definition~\ref{def:op1-one-first-zero-loop-budget} identifies precisely this multiplicity-one situation with the statement that the corridor carries one photon-first-zero transport unit, i.e. \(\Omega_{\mathrm{loop}}=1\). Hence the one-first-zero loop-budget certificate holds. Proposition~\ref{prop:op1-one-first-zero-loop-budget} then yields \(\Delta_{\mathrm{CL}}=0\).
\end{proof}

\begin{corollary}[The live OP~1 burden after r10]
\label{cor:op1-live-burden-after-r10}
After r10, the remaining live burden of OP~1 is no longer to choose a normalization, a pass split, or even an abstract multiplicity slogan. It is to prove a \emph{uniqueness statement on admissible windows}: the retained post-frame corridor supports one and only one primitive photon-first-zero loop compatible with the preserved total-channel load. Any alternative admissible realization would have to appear either as a total-channel distortion or as a second primitive carrier contribution, and therefore lies outside the intended admissible OP~1 corridor class.
\end{corollary}

\begin{remark}[Why r10 is a genuine narrowing and not just another renaming]
This still does not by itself derive the uniqueness statement from first principles; so OP~1 is not yet honestly closed on theorem level. But the queue is now much sharper. Before r10 one could still phrase the task as ``explain multiplicity one'' in a somewhat free way. After r10 the task is pinned to the already existing admissible-window discipline of the present line: one must show that the retained corridor does not admit several distinct primitive first-zero loop realizations once the total-channel profile and no-second-carrier constraints are held fixed. So the next attack on OP~1 should not invent new continuum-limit parameters. It should directly use retained-window rigidity/uniqueness to rule out duplicate admissible loop occupancy on the same corridor.
\end{remark}


\subsection*{19AH$^{(21)}$. Retained-window-rigidity / no-second-carrier bridge for OP~1}
After the admissible-window uniqueness certificate of 19AH$^{(20)}$, the remaining honest task is no longer to invent a new normalization slogan or a new loop-counting parameter. It is to derive uniqueness from structural disciplines that are already native to the present line. The sharpest triad is now visible. First, the retained admissible corridor should be rigid once the retained total-channel profile and optimum data are fixed. Second, the older no-second-carrier discipline should forbid an additional primitive occupancy of that same retained corridor. Third, the strict-upper-gap discipline should forbid a genuinely distinct admissible completion route above the same retained corridor that would duplicate the first-zero budget without altering the present-line status surface. Together these three pressures are exactly what the r10 uniqueness question is asking for in compressed form.

\begin{definition}[Retained-window rigidity certificate for the OP~1 corridor]
\label{def:op1-retained-window-rigidity}
Assume the same-corridor reversal certificate of Definition~\ref{def:op1-same-corridor-reversal}. We say that the profile satisfies the \emph{retained-window rigidity certificate} on the intended admissible windows if every admissible primitive photon-first-zero closed-loop realization on the retained post-frame corridor that preserves the retained total-channel profile, optimum location, and standing admissible equivalence data occupies the same retained corridor slot up to the standing admissible corridor equivalence. Equivalently, there is no nontrivial admissible primitive displacement of the retained first-zero loop inside the same preserved total-channel window.
\end{definition}

\begin{proposition}[Rigidity plus no-second-carrier plus strict upper gap imply OP~1 uniqueness]
\label{prop:op1-rigidity-gap-uniqueness-bridge}
Assume the same-corridor reversal certificate of Definition~\ref{def:op1-same-corridor-reversal}. Assume further that the retained-window rigidity certificate of Definition~\ref{def:op1-retained-window-rigidity} holds on the intended admissible windows, that the no-second-carrier discipline continues to hold at the retained total-channel level, and that the strict upper gap admits no second independent admissible completion route above the same retained corridor on those windows. Then the admissible-window uniqueness certificate of Definition~\ref{def:op1-admissible-window-uniqueness} holds. Consequently,
\[
  \Delta_{\mathrm{CL}} = 0.
\]
In particular, after r11 the live OP~1 burden may be read as the extraction of retained-window rigidity on the post-frame admissible corridor, with the no-second-carrier and strict-upper-gap disciplines supplying the exclusion of duplicate primitive occupancy and duplicate admissible completion.
\end{proposition}

\begin{proof}
Let \(C'\) be any admissible primitive photon-first-zero closed-loop realization on the retained post-frame corridor that preserves the retained total-channel profile. By the retained-window rigidity certificate, \(C'\) occupies the same retained corridor slot as the original loop up to the standing admissible corridor equivalence. If \(C'\) were still a genuinely different primitive realization, then either it would contribute an extra primitive occupancy on that same retained corridor, contradicting the no-second-carrier discipline at total-channel level, or it would represent a second independent admissible completion route above the same corridor while preserving the retained total-channel surface, contradicting the standing strict-upper-gap exclusion on the intended windows. Hence no genuinely distinct admissible primitive realization exists, and the admissible-window uniqueness certificate of Definition~\ref{def:op1-admissible-window-uniqueness} follows. Proposition~\ref{prop:op1-admissible-window-uniqueness} then yields \(\Delta_{\mathrm{CL}}=0\).
\end{proof}

\begin{corollary}[The live OP~1 burden after r11]
\label{cor:op1-live-burden-after-r11}
After r11, OP~1 is no longer honestly read as a continuum-limit mystery in general, nor even as a free uniqueness slogan. The live burden is narrower: extract a retained-window rigidity theorem for the already isolated post-frame admissible corridor. Once that rigidity is in hand, the existing no-second-carrier discipline blocks duplicate primitive corridor occupancy and the strict-upper-gap discipline blocks duplicate admissible completion above the same retained corridor, so the uniqueness required in 19AH$^{(20)}$ follows automatically.
\end{corollary}

\begin{remark}[Why r11 is a real narrowing]
This is still not a final unconditional OP~1 theorem, because the retained-window rigidity certificate itself has not yet been discharged from the active comparison/transport line. But the queue is now much sharper than in r10. The remaining question is no longer ``why uniqueness?'' in the abstract. It is ``why is the retained admissible corridor rigid once the present-line optimum and total-channel data are fixed?'' That is exactly the kind of local structural statement the existing retained-window diagnostics, no-second-carrier discipline, and strict-upper-gap language are already built to support.
\end{remark}


\subsection*{19AH$^{(22)}$. First optimum-pinning comparison-rigidity bridge for retained-window rigidity in OP~1}
After the retained-window-rigidity / no-second-carrier / strict-upper-gap bridge of 19AH$^{(21)}$, the live OP~1 task is sharper still: why should the retained admissible corridor be rigid once the present-line comparison data are fixed? The most honest next compression is to stop treating rigidity as a free geometric slogan and to pin it directly to the already retained comparison surface. On the present line, the retained corridor is not specified by corridor language alone. It is specified together with the preserved total-channel profile, the retained optimum location, and the local comparison/defect data that survive the machine-shadow compression. If those comparison data already determine the admissible slot uniquely inside the retained corridor class, then retained-window rigidity is no longer an extra burden; it is simply the local comparison line read in the right coordinates.

\begin{definition}[Optimum-pinning comparison-rigidity certificate for OP~1]
\label{def:op1-optimum-pinning-comparison-rigidity}
Assume the same-corridor reversal certificate of Definition~\ref{def:op1-same-corridor-reversal}. We say that the profile satisfies the \emph{optimum-pinning comparison-rigidity certificate} on the intended admissible windows if every admissible primitive photon-first-zero closed-loop realization on the retained post-frame corridor that preserves the retained total-channel profile, the retained optimum location, and the standing local comparison/defect data is forced to occupy the same retained corridor slot up to the standing admissible corridor equivalence. Equivalently, within the retained admissible corridor class there is no second primitive loop placement that keeps the same optimum pinning and the same local comparison readout while remaining genuinely distinct.
\end{definition}

\begin{proposition}[Optimum pinning and comparison rigidity imply retained-window rigidity for OP~1]
\label{prop:op1-optimum-pinning-implies-rigidity}
Assume the same-corridor reversal certificate of Definition~\ref{def:op1-same-corridor-reversal}. If the optimum-pinning comparison-rigidity certificate of Definition~\ref{def:op1-optimum-pinning-comparison-rigidity} holds on the intended admissible windows, then the retained-window rigidity certificate of Definition~\ref{def:op1-retained-window-rigidity} holds there as well. Consequently, together with the no-second-carrier and strict-upper-gap disciplines from Proposition~\ref{prop:op1-rigidity-gap-uniqueness-bridge}, one obtains the admissible-window uniqueness certificate of Definition~\ref{def:op1-admissible-window-uniqueness} and therefore
\[
  \Delta_{\mathrm{CL}} = 0.
\]
In particular, after r12 the live OP~1 burden may be read as a local optimum-pinning comparison-rigidity statement on the retained post-frame corridor.
\end{proposition}

\begin{proof}
Take any admissible primitive photon-first-zero closed-loop realization on the retained post-frame corridor that preserves the retained total-channel profile. On the intended admissible windows of the present line, the retained OP~1 corridor is not read independently of the standing comparison surface: the intended loop class is already tracked together with the retained optimum location and the surviving local comparison/defect data of the compressed machine-shadow line. If the optimum-pinning comparison-rigidity certificate holds, then any admissible primitive realization preserving those retained data must occupy the same retained corridor slot up to the standing admissible corridor equivalence. But this is exactly the content of the retained-window rigidity certificate. The rest follows from Proposition~\ref{prop:op1-rigidity-gap-uniqueness-bridge}, which upgrades retained-window rigidity to admissible-window uniqueness under the standing no-second-carrier and strict-upper-gap exclusions, hence yields \(\Delta_{\mathrm{CL}}=0\).
\end{proof}

\begin{corollary}[The live OP~1 burden after r12]
\label{cor:op1-live-burden-after-r12}
After r12, OP~1 is no longer honestly carried as a general rigidity problem. The live burden is narrower: show that the retained optimum and local comparison data already pin the admissible post-frame first-zero corridor slot uniquely. Once that local comparison rigidity is extracted, retained-window rigidity follows, and with the standing no-second-carrier and strict-upper-gap disciplines the OP~1 closure chain completes automatically.
\end{corollary}

\begin{remark}[Why r12 is the first direct attack from the comparison line]
This is still not a full unconditional theorem extracted entirely from the earlier comparison blocks; the optimum-pinning comparison-rigidity certificate itself is now the remaining live statement. But it is a real narrowing. r11 said, in effect, ``prove retained-window rigidity.'' r12 says something more local and more testable: ``prove that no genuinely distinct primitive loop placement preserves the retained optimum and the surviving comparison/defect readout on the same admissible corridor.'' That is the first point at which the residual OP~1 burden is read directly as a comparison-line pinning problem rather than as a free rigidity slogan.
\end{remark}

\subsection*{19AI. Bundled overlay-completion capstone for the machine-theoretic line}
After the machine-theoretic line has reached full machine-theoretic closure, present-line status readout, and window-equivalence, the remaining overlay refinements should no longer be read as an invitation to keep appending parallel adjectives indefinitely. Their real purpose is to certify, in a bundled way, that the machine layer is the right kind of organization of the already certified machine-bearing windows: it adds no foreign admissible content, forgets none of the operative distinctions relevant for the machine hierarchy, tracks the intended statuses exactly, and returns the same machine-bearing layer up to the standing admissible equivalences. The correct freeze-level conclusion is therefore not a new ladder of further micro-properties, but a single capstone statement that the machine overlay is a conservative, faithful, exact, relatively complete, adequate, reconstructive, closed, idempotent, minimal, canonical, invariant, natural, universally characterized, rigid, terminal, and equivalence-level organization of the intended certified machine-bearing windows.

\begin{definition}[Bundled machine-overlay completion certificate on the compressed computational line]
Fix a chosen sequential reference core \(\mathcal R\) and an intended admissible machine-bearing window class \(\mathcal W^{\mathrm{mach}}\). We say that the compressed computational line satisfies the \emph{bundled machine-overlay completion certificate} on \(\mathcal W^{\mathrm{mach}}\) if, once the full machine-theoretic closure surface \(\mathfrak M^{\mathrm{full}}(\mathcal R)=1\), the present-line status surface \(\mathfrak P^{\mathrm{mach}}(\mathcal R)\), and the window-equivalence/equivalence-surface layer hold on \(\mathcal W^{\mathrm{mach}}\), the resulting machine-theoretic overlay is conservative, faithful, exact, relatively complete, adequate, reconstructive, closed, idempotent, minimal, canonical, invariant, natural, universally characterized, rigid, terminal, and equivalent to the certified machine-bearing structural content, all up to the standing admissible window-equivalence and overlay-equivalence. Equivalently, the intended admissible machine-bearing windows and the canonical machine-overlay objects form two fully compatible organizations of the same certified machine-bearing layer.
\end{definition}

\begin{proposition}[Bundled overlay-completion theorem for the machine-theoretic line]
On the intended admissible machine-bearing windows of the present compressed computational line, the above bundled machine-overlay completion certificate holds. More precisely, the machine-theoretic line of Sections~19J--19AH together with the diagnostic continuation block 19AH$\prime$--19AH$^{(22)}$ yields the indispensable ladder from local operational semantics to the present-line conditional/unconditional machine status plus the first disciplined re-reading of the retained machine shadow, while the subsequent overlay-completion work can be read in bundled form as showing that this machine layer is conservative, faithful, exact, relatively complete, adequate, reconstructive, closed, idempotent, minimal, canonical, invariant, natural, universally characterized, rigid, terminal, and equivalent to the intended certified machine-bearing windows. Consequently, for the present Companion purpose the machine-theoretic overlay may be cited as a closed reversible machine-level organization of the certified machine-bearing structural content, without needing to unfold each of these overlay attributes separately unless a later revision introduces a genuinely new machine-theoretic distinction.
\end{proposition}

\begin{corollary}[What the bundled overlay-completion capstone already buys]
The machine line now reaches a true freeze-level stopping point. For subsequent work it is enough to cite the indispensable machine ladder together with the bundled overlay-completion capstone, rather than to continue lengthening the line by parallel closure adjectives. In particular, the present line can now be read through the three decisive bundled surfaces --- the full machine-theoretic closure surface, the present-line status surface, and the bundled overlay-completion capstone --- while the inserted channel-split diagnostics remain the preferred local reading tools whenever carrier and observer-side burdens need to be separated.
\end{corollary}

